\subsection{The canonical injective resolution}

Let F be the A-module \(\mathrm{Hom}_{\mathbf{Z}}(\mathbf{A},\mathbf{Q} / \mathbf{Z})\) ; for every A-module M, we set \(\mathrm{I}^0 (\mathbf{M}) = \mathrm{F}^{\mathrm{Hom}_{\mathbf{A}}(\mathbf{M},\mathbf{F})}\) and we denote by \(e_{\mathbf{M}}:\mathbf{M}\to \mathrm{I}^{0}(\mathbf{M})\) the homomorphism which maps \(m\in \mathbf{M}\) to the family \((\varphi (m))_{\varphi \in \mathrm{Hom}_{\mathbf{A}}(\mathbf{M},\mathbf{F})}\) . By X, p.~19, cor. 2, \(\mathrm{I}^0 (\mathbf{M})\) is an injective A-module and \(e_{\mathbf{M}}\) is injective. Set \(\mathbf{K}^0 (\mathbf{M}) = \mathbf{Coker}\) \(e_{\mathbf{M}}\) and denote by \(q_{\mathbf{M}}:\mathrm{I}^0 (\mathbf{M})\to \mathrm{K}^0 (\mathbf{M})\) the canonical projection. We thus have an exact sequence

\[
0 \longrightarrow \mathrm{M} \stackrel {e _ {\mathrm{M}}} {\longrightarrow} \mathrm{I} ^ {0} (\mathrm{M}) \stackrel {q _ {\mathrm{M}}} {\longrightarrow} \mathrm{K} ^ {0} (\mathrm{M}) \longrightarrow 0.
\]

We define a graded A-module I(M) by setting \(\mathbf{I}^n (\mathbf{M}) = 0\) for \(n < 0\) and, by induction on the integer \(n > 0\) ,

\[
\mathrm{I} ^ {n} (\mathbf {M}) = \mathrm{I} ^ {0} \left(\mathrm{K} ^ {n - 1} (\mathbf {M})\right), \quad \mathrm{K} ^ {n} (\mathbf {M}) = \mathrm{K} ^ {0} \left(\mathrm{K} ^ {n + 1} (\mathbf {M})\right).\tag{9}
\]

We define A-homomorphisms \(\delta_{\mathbf{M}}^{n}:\mathrm{I}^{n}(\mathbf{M})\to \mathrm{I}^{n + 1}(\mathbf{M})\) by

\[
\left\{ \begin{array}{l l} \delta_ {\mathbf {M}} ^ {n} = 0, & n <   0, \\ \delta_ {\mathbf {M}} ^ {0} = e _ {\mathbf {K} ^ {0} (\mathbf {M})} \circ q _ {\mathbf {M}}, \\ \delta_ {\mathbf {M}} ^ {n} = e _ {\mathbf {K} ^ {n} (\mathbf {M})} \circ q _ {\mathbf {K} ^ {n - 1} (\mathbf {M})}, & n > 0. \end{array} \right.\tag{10}
\]

By construction, we have an exact sequence

\[
0 \longrightarrow \mathrm{M} \xrightarrow {e _ {\mathrm{M}}} \mathrm{I} ^ {0} (\mathrm{M}) \xrightarrow {\delta_ {\mathrm{M}} ^ {0}} \dots \longrightarrow \mathrm{I} ^ {n} (\mathrm{M}) \xrightarrow {\delta_ {\mathrm{M}} ^ {n}} \mathrm{I} ^ {n + 1} (\mathrm{M}) \longrightarrow \dots ,
\]

so that, if one extends \(e_{M}\) as a morphism of complexes

\[
e _ {\mathbf {M}}: \mathbf {M} \rightarrow (\mathrm{I} (\mathbf {M}), \delta_ {\mathbf {M}}),
\]

we obtain an injective resolution of M, called the canonical injective resolution of M. Let \(f: \mathbf{M} \to \mathbf{N}\) be a homomorphism of A-modules. Denote by \(\mathrm{I}^{0}(f)\) the homomorphism from \(\mathrm{I}^{0}(\mathbf{M}) = \mathrm{F}^{\mathrm{Hom}_{\mathbf{A}}(\mathbf{M},\mathbf{F})}\) to \(\mathrm{I}^{0}(\mathbf{N}) = \mathrm{F}^{\mathrm{Hom}_{\mathbf{A}}(\mathbf{N},\mathbf{F})}\) which maps the family \((x_{\varphi})_{\varphi \in \mathrm{Hom}_{\mathbf{A}}(\mathbf{M},\mathbf{F})}\) to the family \((x_{\psi \circ f})_{\psi \in \mathrm{Hom}_{\mathbf{A}}(\mathbf{N},\mathbf{F})}\) . We have:

\[
\mathrm{I} ^ {0} (f) \circ e _ {\mathbf {M}} = e _ {\mathbf {N}} \circ f.\tag{11}
\]

Consequently, \(\mathrm{I}^0 (f)\) induces a homomorphism \(\mathbf{K}^0 (f):\mathbf{K}^0 (\mathbf{M})\to \mathbf{K}_-^0 (\mathbf{N})\) and we have

\[
\mathrm{K} ^ {0} (f) \circ q _ {\mathbf {M}} = q _ {\mathbf {N}} \circ \mathrm{K} ^ {0} (f).\tag{12}
\]

Set \(\mathrm{I}^n(f) = 0\) for \(n < 0\) and let us define, by induction on the integer \(n > 0\), homomorphisms \(\mathrm{I}^n(f): \mathrm{I}^n(\mathbf{M}) \to \mathrm{I}^n(\mathbf{N})\) and \(\mathrm{K}^n(f): \mathrm{K}^n(\mathbf{M}) \to \mathrm{K}^n(\mathbf{N})\) by:

\[
\left\{ \begin{array}{l} \mathrm{I} ^ {n} (f) = \mathrm{I} ^ {0} (\mathrm{K} ^ {n - 1} (f)) \\ \mathrm{K} ^ {n} (f) = \mathrm{K} ^ {0} (\mathrm{K} ^ {n - 1} (f)). \end{array} \right.\tag{13}
\]

PROPOSITION 5. --- I(f) : I(M) → I(N) is a morphism of complexes of A-modules; we have I(f) ∘ e\_M = e\_N ∘ f.

This is proved in a manner analogous to prop. 4.

One has

\[
\mathrm{I} (1 _ {\mathbf {M}}) = 1 _ {\mathrm{I} (\mathbf {M})}\tag{14}
\]

and for every homomorphism \(g: N \rightarrow P\) of A-modules

\[
\mathrm{I} (g \circ f) = \mathrm{I} (g) \circ \mathrm{I} (f).\tag{15}
\]

Remark. --- If \(f, g \in \operatorname{Hom}_{\mathbf{A}}(\mathbf{M}, \mathbf{N})\) , we do not have \(\mathrm{I}(f + g) = \mathrm{I}(f) + \mathrm{I}(g)\) . However, these two morphisms are homotopic by X, p.~49, prop. 3 bis.

If M is a right A-module, we set \(\mathrm{I}(\mathbf{M}) = \mathrm{I}(\mathbf{M}^{\circ})^{\circ}\) ; it is called the canonical injective resolution of M and we have

\[
\mathrm{I} (\mathbf {M} ^ {\circ}) = \mathrm{I} (\mathbf {M}) ^ {\circ}.
\]

\subsection{Finitely generated resolutions}

It follows in particular from the two preceding sections that every A-module has injective resolutions, free resolutions (hence also projective or flat resolutions). In certain cases, one can be more precise:

Suppose A is left Noetherian and let M be an A-module. Let us construct recursively sequences \((\mathbf{L}_{n})_{n\geqslant0},(\mathbf{Z}_{n})_{n\geqslant0},(d_{n})_{n\geqslant1}\) where, for all \(n\geqslant0\) , \(L_{n}\) is a finitely generated free A-module, \(Z_{n}\) a submodule of \(L_{n}\) and \(d_{n+1}:L_{n+1}\to L_{n}\) a homomorphism. For this, choose a finite generating family \((m_{i})_{i\in I_{0}}\) of M, set \(L_{0}=A^{(I_{0})}\) , define \(p:L_{0}\to M\) by \(p(e_{i})=m_{i}\) and set \(Z_{0}=\operatorname{Ker}(p)\) . For \(n\geqslant0\) , the modules \(L_{n}\) and \(Z_{n}\) having been constructed, \(Z_{n}\) is of finite type since contained in \(L_{n}\) ; let us choose a finite generating family \((x_{n,i})_{i\in I_{n+1}}\) of \(Z_{n}\) ; set \(L_{n+1}=A^{(I_{n+1})}\) , define \(d_{n+1}\) by \(d_{n+1}(e_{i})=x_{n,i}\) and set \(Z_{n+1}=\operatorname{Ker}(d_{n+1})\) . By construction, we have an exact sequence

\[
\dots \longrightarrow \mathrm{L} _ {n + 1} \xrightarrow {d _ {n + 1}} \mathrm{L} _ {n} \longrightarrow \dots \longrightarrow \mathrm{L} _ {0} \xrightarrow {p} \mathrm{M} \longrightarrow 0,
\]

whence:

PROPOSITION 6. --- When A is left Noetherian, every finitely generated A-module M has a free resolution \(p: \mathbf{L} \to \mathbf{M}\) such that \(\mathbf{L}_n\) is finitely generated for every integer \(n\) . More generally:

PROPOSITION 7. --- Let C be an A-complex and let \(a \in \mathbf{Z}\) be such that \(\mathrm{H}_{n}(\mathbf{C}) = 0\) for \(n < a\) . a) There exists a free A-complex \(\mathbf{L}\) such that \(\mathbf{L}_{n} = 0\) for \(n < a\) and a homologism \(f: \mathbf{L} \to \mathbf{C}\) .

\begin{enumerate}
\def\labelenumi{\alph{enumi})}
\setcounter{enumi}{1}
\tightlist
\item
  Suppose A is left Noetherian and the A-modules \(\mathrm{H}_{n}(\mathbf{C})\) , \(n \in Z\) , are finitely generated. There exists a free A-complex L such that \(L_{n} = 0\) for n \textless{} a and that \(L_{n}\) be a finitely generated A-module for every n, and a homologism \(f: L \to C\) .
\end{enumerate}

Let C' be the subcomplex of C such that \(C_{n}^{\prime}=C_{n}\) for n\textgreater a, \(\mathbf{C}_{a}^{\prime}=\mathbf{Z}_{a}(\mathbf{C})\) , \(C_{n}^{\prime}=0\) for n \textless{} a; then the canonical injection from C' into C is a homologism. Replacing C by \(C'\) , we may therefore assume that \(C_{n}=0\) for n\textless a. The statement then follows from the iterated application of the following lemma, for r=a, \(a+1,\ldots\) :

Lemma 3. --- Let C be a complex and \(r \in \mathbf{Z}\) . There exists a complex \(\mathbf{C}'\) and a homologism \(f: \mathbf{C}' \to \mathbf{C}\) such that \(f_n: \mathbf{C}_n' \to \mathbf{C}_n\) be an isomorphism for \(n < r\) and that \(\mathbf{C}_r'\) be a free A-module. If A is Noetherian and the A-modules \(\mathrm{H}_r(\mathbf{C})\) and \(\mathbf{C}_{r-1}\) are finitely generated, one can impose that \(\mathbf{C}_r'\) be finitely generated.

\begin{enumerate}
\def\labelenumi{\alph{enumi})}
\item
  First let \(h: \mathbf{M} \to \mathbf{C}_r\) be a homomorphism of A-modules; denote by \(d = (d_n)\) the differential of C. Let N be the submodule of \(\mathbf{M} \times \mathbf{C}_{r+1}\) formed by the pairs \((m, x)\) such that \(h(m) = d_{r+1}(x)\) ; we define a complex \((\mathbf{C}', d')\) by \(\mathbf{C}_n' = \mathbf{C}_n\) for \(n \neq r, r+1, \mathbf{C}_r' = \mathbf{M}, \mathbf{C}_{r+1}' = \mathbf{N}, d_n' = d_n\) for \(n \neq r, r+1, r+2, d_r' = d_r \circ h, d_{r+1}'(m, x) = m\) for \((m, x) \in \mathbf{N}\) and \(d_{r+2}'(y) = (0, d_{r+2}(y))\) for \(y \in \mathbf{C}_{r+2}\) . Consider also the morphism of complexes \(f: \mathbf{C}' \to \mathbf{C}\) such that \(f_n = 1_{\mathbf{C}_n}\) for \(n \neq r, r+1, f_r = h, f_{r+1}(m, x) = x\) .
\item
  The complex Ker f is zero in degree \(\neq r, r + 1\) and the differential \(d_{r+1}^{\prime}\) induces an isomorphism from \(\operatorname{Ker} f_{r+1}\) onto \(\operatorname{Ker} f_r\) , hence \(\operatorname{H}(\operatorname{Ker} f) = 0\) .
\item
  When the composite map \(\mathbf{M} \xrightarrow{\hbar} \mathbf{C}_r \to \mathbf{C}_r / \mathbf{B}_r(\mathbf{C})\) is surjective, one sees similarly that \(\mathrm{H}(\mathrm{Coker}f) = 0\) , and \(f\) is then a homologism (X, p.~31, cor. 2).
\item
  When A is assumed to be Noetherian and the A-modules \(H_{r}(C)\) and \(C_{r-1}\) are finitely generated, then \(C_{r}/B_{r}(C)\) is finitely generated, by virtue of the exact sequence (X, p.~25)
\end{enumerate}

\[
0 \to \mathrm{H} _ {r} (\mathrm{C}) \to \mathrm{C} _ {r} / \mathrm{B} _ {r} (\mathrm{C}) \to \mathrm{C} _ {r - 1};
\]

there then exists a finitely generated free A-module M and a homomorphism \(h: \mathbf{M} \to \mathbf{C}_{r}\) such that the condition of \(c)\) be satisfied; in the general case, there exists a free module M and a surjective homomorphism \(h: \mathbf{M} \to \mathbf{C}_{r}\) . This completes the proof.

\subsection{Minimal projective resolutions}

Let M be an A-module and let

\[
\dots \longrightarrow \mathrm{P} _ {n} \xrightarrow {d _ {n}} \mathrm{P} _ {n - 1} \longrightarrow \dots \longrightarrow \mathrm{P} _ {0} \xrightarrow {d _ {0}} \mathrm{M} \longrightarrow 0\tag{P}
\]

a resolution of M. We say that (P) is a minimal projective resolution if, for every \(n \geqslant 0\) , the homomorphism \(\delta_{n}: P_{n} \to \operatorname{Im}(d_{n})\) induced by \(d_{n}\) is a projective cover (VIII, § 8, no. 5).

PROPOSITION 8. — Let M be an A-module, and let P and P' be two minimal projective resolutions of M, with f: P → P' a morphism of resolutions. Then f is an isomorphism. In particular, any two minimal projective resolutions of M are isomorphic.

Set \(\tilde{\mathbf{P}}_{n} = \mathbf{P}_{n}\) for \(n \neq -1\) and \(\tilde{\mathbf{P}}_{-1} = \mathbf{M}\) ; let us define similarly \(\tilde{\mathbf{P}}_{n}^{\prime}\) and set \(f_{-1} = 1_{\mathbf{M}}\) . Let us prove by induction starting from \(-1\) that \(f_{n}: \tilde{\mathbf{P}}_{n} \to \tilde{\mathbf{P}}_{n}^{\prime}\) is an isomorphism for all \(n\) . This is evident for \(n = -1\) ; suppose that \(f_{n}\) and \(f_{n-1}\) are isomorphisms. It follows from the commutativity of the diagram:

\[
\begin{array}{c c c} \mathbf {P} _ {n} & \xrightarrow {d _ {n}} & \mathbf {P} _ {n - 1} \\ \Big \downarrow f _ {n} & & \Big \downarrow f _ {n - 1} \\ \mathbf {P} _ {n} ^ {\prime} & \xrightarrow {d _ {n} ^ {\prime}} & \mathbf {P} _ {n - 1} ^ {\prime} \end{array}
\]

that \(f_{n}\) induces an isomorphism \(g_{n}\) of \(\operatorname{Ker} d_{n}\) onto \(\operatorname{Ker} d_{n}^{\prime}\) . It then follows from the commutativity of the diagram:

\[
\begin{array}{c c c} \mathrm{P} _ {n + 1} & \xrightarrow {\delta_ {n + 1}} & \text {Ker} d _ {n} \\ \Big \downarrow f _ {n + 1} & & \Big \downarrow g _ {n} \\ \mathrm{P} _ {n + 1} ^ {\prime} & \xrightarrow {\delta_ {n + 1} ^ {\prime}} & \text {Ker} d _ {n} ^ {\prime} \end{array}
\]

and from VIII, loc. cit., that \(f_{n+1}\) is an isomorphism.

COROLLARY. --- Let M be an A-module, and P and P' two projective resolutions of M; suppose that P is minimal. Let f : P → P' and g : P' → P be two morphisms of resolutions. Then f is injective, g is surjective, and P' is the direct sum of the subcomplexes Im f and Ker g. Moreover Ker g has zero homology.

Indeed, \(\alpha = g \circ f\) is an automorphism of P (prop. 8). Set \(\tilde{f} = f \circ \alpha^{-1}\) . We have

\[
\operatorname{Im} \tilde {f} = \operatorname{Im} f \quad \text { and } \quad g \circ \tilde {f} = 1 _ {\mathrm{P}},
\]

which shows that \(\mathbf{P}' = \operatorname{Im} \tilde{f} \oplus \operatorname{Ker} g\) . Since the sequence

\[
0 \rightarrow \operatorname{Ker} g \rightarrow \mathrm{P} ^ {\prime} \stackrel {g} {\rightarrow} \mathrm{P} \rightarrow 0
\]

is exact and \(g\) is a homologism, \(\operatorname{Ker} g\) has zero homology.

PROPOSITION 9. --- Let M be an A-module and (P, p) a projective resolution of M. Let r denote the radical of A. Suppose either that \(P_{n}\) is a finitely generated A-module for all n, or that r is nilpotent. Then for (P, p) to be minimal it is necessary and sufficient that the complex \((\mathbf{A}/\mathbf{r}) \otimes_{\mathbf{A}} \mathbf{P}\) have zero differential, in other words that

\[
d _ {n + 1} (\mathbf {P} _ {n + 1}) \subset \mathrm{r} \mathbf {P} _ {n} \quad \text { for all } n \geqslant 0.
\]

Suppose that (P, p) is minimal. By VIII, loc. cit., the homomorphism

\[
1 \otimes \delta_ {n}: (\mathrm{A} / \mathrm{r}) \otimes_ {\mathrm{A}} \mathrm{P} _ {n} \rightarrow (\mathrm{A} / \mathrm{r}) \otimes_ {\mathrm{A}} \operatorname{Im} d _ {n}
\]

is an isomorphism. It then follows from the exact sequence

\[
0 \longrightarrow \operatorname{Im} d _ {n + 1} \xrightarrow {j _ {n}} \mathrm{P} _ {n} \xrightarrow {\delta_ {n}} \operatorname{Im} d _ {n} \longrightarrow 0
\]

that the homomorphism \(1 \otimes j_n : (\mathbf{A} / \mathfrak{r}) \otimes_{\mathbf{A}} \operatorname{Im} d_{n+1} \to (\mathbf{A} / \mathfrak{r}) \otimes_{\mathbf{A}} \mathrm{P}_n\) is zero; since \(d_{n+1} = j_n \circ \delta_{n+1}\) , we deduce that \(1_{\mathbf{A} / \mathfrak{r}} \otimes d_{n+1} = 0\) for \(n \geqslant 0\) .

Conversely, suppose that for all \(n \geqslant 1\) , \(1 \otimes d_n\) is zero, in other words that \(\operatorname{Im} d_n = \operatorname{Ker} d_{n-1}\) is contained in \(\mathbf{rP}_{n-1}\) . Since \(\delta_{n-1}\) is surjective, it follows from VIII, loc. cit. that \(\delta_{n-1}\) is a projective cover for \(n \geqslant 1\) , hence that \((\mathbf{P}, p)\) is minimal.

PROPOSITION 10. --- Suppose that A is a left Noetherian local ring, and let M be a finitely generated A-module. Then M possesses a minimal resolution (P, p); for every \(n \geqslant 0\) , \(\mathbf{P}_{n}\) is a finitely generated free module.

Indeed, in the construction made in n° 5 (p.~53), one may, by virtue of VIII, loc. cit., take for \((\mathrm{L}_{0}, p)\) a projective cover of M, and for \(L_{n+1}\) a projective cover of Ker \(d_{n}\) . The resolution thus obtained is then minimal.

Remarks. --- 1) Let m denote the maximal ideal of A and set \(k = \mathbf{A}/\mathfrak{m}\) . Let P be a minimal projective resolution of M, and set \(b_{n} = \dim_{k}(k \otimes_{\mathbf{A}} \mathbf{P}_{n})\) . Then \(\mathbf{P}_{n}\) is a free A-module of rank \(b_{n}\) . It follows from the corollary of prop. 8 that for any other projective resolution \(\mathbf{P}'\) of M, one has \(\dim_{k}(k \otimes_{\mathbf{A}} \mathbf{P}_{n}') \geqslant b_{n}\) , and that equality holds if and only if \(\mathbf{P}'\) is minimal.

\begin{enumerate}
\def\labelenumi{\arabic{enumi})}
\setcounter{enumi}{1}
\tightlist
\item
  By prop. 9, \(b_n\) is the dimension over \(k\) of \(H_n(k \otimes_A P)\) , * in other words of \(\text{Tor}_n^A(k, M)\) . It is also the dimension over \(k\) of \(\text{Ext}_A^n(M, k)\) (cf.~X, p.~103, remark 3) *.
\end{enumerate}

\subsection{Graded resolutions}

In this section, one supposes that the ring A is equipped with a grading \((\mathbf{A}_n)_{n\in \mathbf{Z}}\) , such that \(\mathbf{A}_n = 0\) for \(n < 0\) . One says that a graded A-module M is bounded below if \(\mathbf{M}_n = 0\) for n sufficiently small; every finitely generated graded A-module is bounded below.

PROPOSITION 11. --- If M is a graded A-module bounded below (resp. if M is a finitely generated graded A-module and if A is left Noetherian), there exists an exact sequence of graded A-modules, unbounded to the left,

\[
\dots \longrightarrow \mathrm{L} _ {n} \xrightarrow {d _ {n}} \mathrm{L} _ {n - 1} \longrightarrow \dots \longrightarrow \mathrm{L} _ {1} \xrightarrow {d _ {1}} \mathrm{L} _ {0} \xrightarrow {d _ {0}} \mathrm{M} \longrightarrow 0
\]

where the \(L_{i}\) are graded free and bounded below (resp. graded free and finitely generated), and where the \(d_{i}\) are graded homomorphisms of degree 0.

If N is a graded A-module bounded below (resp. and finitely generated over Noetherian A) there exists a graded free A-module bounded below (resp. and finitely generated) L and a surjective graded homomorphism L → N (II, p.~167, remark 3).

That being so, suppose given an exact sequence of graded A-modules and graded homomorphisms of degree 0

\[
\mathrm{L} _ {n} \xrightarrow {d _ {n}} \mathrm{L} _ {n - 1} \longrightarrow \dots \longrightarrow \mathrm{L} _ {0} \xrightarrow {d _ {0}} \mathrm{M} \longrightarrow 0,
\]

where the \(L_{i}, i = 0, \ldots, n\) , are graded free and bounded below (resp. graded free and finitely generated). Then N = Ker \(d_{n}\) is bounded below (resp. finitely generated); there therefore exists a graded free A-module bounded below (resp. graded free and finitely generated) \(L_{n+1}\) and a graded homomorphism \(d_{n+1}: L_{n+1} \to L_{n}\) of degree 0, such that Im \(d_{n+1} = N\) ; the sequence

\[
\mathrm{L} _ {n + 1} \xrightarrow {d _ {n + 1}} \mathrm{L} _ {n} \xrightarrow {d _ {n}} \mathrm{L} _ {n - 1} \longrightarrow \dots \longrightarrow \mathrm{L} _ {0} \xrightarrow {d _ {0}} \mathrm{M} \longrightarrow 0
\]

is then exact. The proposition then follows from the preceding by induction on n.

\subsection{The standard resolution}

In this section, one supposes that the ring A is an (associative and unital) algebra over a commutative ring \(k\) . For \(n \geqslant 0\) we denote by \(\mathbf{B}_n\) the tensor product over \(k\) of \((n + 2)\) modules equal to A. We regard it as an (A, A)-bimodule by endowing it with the left (resp. right) A-module structure induced by the left (resp. right) A-module structure of the first (resp. last) factor of the tensor product.

For \(n \geqslant 1\) , we define bimodule homomorphisms \(d_n^i\) (for \(0 \leqslant i \leqslant n\) ) and \(d_n\) from \(B_n\) to \(B_{n-1}\) , by the formulas:

\[
d _ {n} ^ {i} (x _ {0} \otimes \ldots \otimes x _ {n + 1}) = x _ {0} \otimes \ldots \otimes x _ {i} x _ {i + 1} \otimes \ldots \otimes x _ {n + 1}, \quad 0 \leqslant i \leqslant n,
\]

\[
d _ {n} = \sum_ {i = 0} ^ {n} (- 1) ^ {i} d _ {n} ^ {i}.
\]

It is clear that

\[
d _ {n - 1} ^ {i} \circ d _ {n} ^ {j} = d _ {n - 1} ^ {j - 1} \circ d _ {n} ^ {i} \quad \text { for } \quad i <   j
\]

and consequently

\[
d _ {n - 1} \circ d _ {n} = \sum_ {0 \leqslant i <   j \leqslant n} (- 1) ^ {i + j} d _ {n - 1} ^ {i} \circ d _ {n} ^ {j} + \sum_ {0 \leqslant j \leqslant i \leqslant n - 1} (- 1) ^ {i + j} d _ {n - 1} ^ {i} \circ d _ {n} ^ {j} = 0.
\]

Therefore, if we set \(\mathbf{B}_n = 0\) for \(n < 0\) and \(d_n = 0\) for \(n \leqslant 0\) , the sequence \((\mathbf{B}_n, d_n)\) defines a complex of (A, A)-bimodules (X, p.~43), which will be denoted \(\mathbf{B}(\mathbf{A})\) . For every left A-module \(\mathbf{M}\) we denote by \(\mathbf{B}(\mathbf{A}, \mathbf{M})\) the complex formed by the \(\mathbf{B}_n \otimes_{\mathbf{A}} \mathbf{M}\) and the \(d_n \otimes 1_{\mathbf{M}}\) , in other words the tensor product complex \(\mathbf{B}(\mathbf{A}) \otimes_{\mathbf{A}} \mathbf{M}_{*}\) ; it is a complex of left A-modules.

We define an A-linear map \(\varepsilon_{\mathbf{M}}\) from \(\mathbf{B}_{0}(\mathbf{A},\mathbf{M}) = (\mathbf{A}\otimes_{k}\mathbf{A})\otimes_{\mathbf{A}}\mathbf{M}\) to \(\mathbf{M}\) by the formula \(\varepsilon_{\mathbf{M}}(a\otimes b\otimes m) = abm\) for \(a,b\in \mathbf{A},m\in \mathbf{M}\) . We have \(\varepsilon_{\mathbf{M}}\circ d_1 = 0\) , so that the graded homomorphism \(\overline{\varepsilon}_{\mathbf{M}}:\mathbf{B}(\mathbf{A},\mathbf{M})\to \mathbf{M}\) , which coincides with \(\varepsilon_{\mathbf{M}}\) in degree 0, is a morphism of complexes of A-modules.

PROPOSITION 12. --- The map \(\overline{\varepsilon}_{\mathbf{M}}: \mathbf{B}(\mathbf{A}, \mathbf{M}) \to \mathbf{M}\) is a homotopism of complexes of \(k\) -modules. In particular, the complex \(\mathbf{B}(\mathbf{A}, \mathbf{M})\) is split over \(k\) , and \((\mathbf{B}(\mathbf{A}, \mathbf{M}), \overline{\varepsilon}_{\mathbf{M}})\) is a left resolution of the A-module M.

For \(n \geqslant 0\) , let us define a \(k\) -linear map \(s_n: \mathbf{B}_n \to \mathbf{B}_{n+1}\) by the formula:

\[
s _ {n} \left(x _ {0} \otimes \dots \otimes x _ {n + 1}\right) = 1 \otimes x _ {0} \otimes \dots \otimes x _ {n + 1} \quad \text { for } \quad x _ {0}, \dots , x _ {n + 1} \in A.
\]

It is a homomorphism of right A-modules, which satisfies the identities:

\[
d _ {n + 1} ^ {i} \circ s _ {n} = s _ {n - 1} \circ d _ {n} ^ {i - 1} \quad \mathrm{for} \quad n \geqslant 1, \quad 1 \leqslant i \leqslant n + 1,
\]

\[
d _ {n + 1} ^ {0} \circ s _ {n} = 1 _ {\mathrm{B} _ {n}} \quad \text { for } \quad n \geqslant 1,
\]

and consequently we have

\[
d _ {n + 1} \circ s _ {n} + s _ {n - 1} \circ d _ {n} = 1 _ {\mathrm{B} _ {n}} \quad \text { for } \quad n \geqslant 1.\tag{16}
\]

Moreover, we have

\[
d _ {1} \circ s _ {0} (x _ {0} \otimes x _ {1}) = x _ {0} \otimes x _ {1} - 1 \otimes x _ {0} x _ {1} \quad \text { for } \quad x _ {0}, x _ {1} \in A.\tag{17}
\]

Let us denote by \(\eta: \mathbf{A} \to \mathbf{A} \otimes_k \mathbf{A}\) the map defined by \(\eta(a) = 1 \otimes a\) , and \(\overline{\eta}: \mathbf{A} \to \mathbf{B}(\mathbf{A})\) the morphism of complexes which coincides with \(\eta\) in degree 0. It is clear that \(\overline{\varepsilon}_{\mathbf{A}} \circ \overline{\eta} = 1_{\mathbf{A}}\) ; formulas (16) and (17) show that \(d \circ s + s \circ d = 1_{\mathrm{B(A)}} - \overline{\eta} \circ \overline{\varepsilon}_{\mathbf{A}}\) . Setting \(\overline{\eta}_{\mathrm{M}} = \overline{\eta} \otimes 1_{\mathrm{M}}\) , \(d_{\mathrm{M}} = d \otimes 1_{\mathrm{M}}\) and \(s_{\mathrm{M}} = s \otimes 1_{\mathrm{M}}\) , we deduce that \(\overline{\varepsilon}_{\mathrm{M}} \circ \overline{\eta}_{\mathrm{M}} = 1_{\mathrm{M}}\) and \(d_{\mathrm{M}} \circ s_{\mathrm{M}} + s_{\mathrm{M}} \circ d_{\mathrm{M}} = 1_{\mathrm{B(A,M)}} - \overline{\eta}_{\mathrm{M}} \circ \overline{\varepsilon}_{\mathrm{M}}\) . In other words, (X, p.~33, def. 5), \(\overline{\varepsilon}_{\mathrm{M}}\) is a homotopism of complexes of \(k\) -modules. The other assertions of the proposition follow immediately.

DEFINITION 3. --- The left resolution (B(A, M), \(\overline{\varepsilon}_{\mathrm{M}}\) ) of M is called the standard resolution of the A-module M.

If A and M are projective (resp. free, resp. flat) k-modules, the standard resolution B(A, M) is a projective (resp. free, resp. flat) resolution of M.

\subsection{Resolutions and Grothendieck groups}

If \(\mathcal{C}\) is a set of classes of A-modules, we shall say that a left resolution \((\mathbf{P}, p)\) is bounded of type \(\mathcal{C}\) if the complex \(\mathbf{P}\) is bounded of type \(\mathcal{C}\) (X, p.~41).

THEOREM 1. --- Let \(\mathcal{C}_{0}\) and \(\mathcal{C}\) be two additive and left exact classes of A-modules such that \(\mathcal{C}_{0} \subset \mathcal{C}\) and that every A-module of type \(\mathcal{C}\) has a bounded left resolution of type \(\mathcal{C}_{0}\) . Then the homomorphism \(\alpha : \mathrm{K}(\mathcal{C}_{0}) \to \mathrm{K}(\mathcal{C})\) deduced from the inclusion of \(\mathcal{C}_{0}\) into \(\mathcal{C}\) is bijective; if M is an A-module of type \(\mathcal{C}\) and P a bounded left resolution of M of type \(\mathcal{C}_{0}\) , we have \(\alpha^{-1}([M]_{\mathcal{C}}) = \chi_{\mathcal{C}_{0}}(P)\) (X, p.~41, example 6).

Lemma 4. --- Let \(f: \mathbf{M}' \to \mathbf{M}\) be a homomorphism of A-modules of type \(\mathcal{C}\) , and \(p: \mathbf{P} \to \mathbf{M}\) a left resolution of \(\mathbf{P}\) bounded of type \(\mathcal{C}_0\) There exists a left resolution \(p': \mathbf{P}' \to \mathbf{M}'\) bounded of type \(\mathcal{C}_0\) and a morphism of complexes \(u: \mathbf{P}' \to \mathbf{P}\) such that \(p \circ u = f \circ p'\) .

Let us reason by induction on the length n of P, the assertion being trivial when this length is \textless{} 0. Consider the mapping \(g: \mathbf{M}' \times \mathbf{P}_0 \to \mathbf{M}\) such that

\[
g (x, y) = f (x) - p _ {0} (y) \quad \text { for } \quad x \in \mathbf {M} ^ {\prime}, y \in \mathbf {P} _ {0},
\]

and its kernel K; the A-module K is of type C since g is surjective and \(M' \times P_{0}\) and M are of type C. Let \(h : P_{0}' \to K\) be a surjective homomorphism where \(P_{0}'\) is of type \(C_{0}\) ; denote by \(p_{0}' : P_{0}' \to M'\) (resp. \(u_{0} : P_{0}' \to P_{0}\) ) the composite homomorphism of h and the projection \(K \to M\) (resp. \(K \to P_{0}\) ); the homomorphism \(p_{0}'\) is surjective and we have a commutative diagram

\[
\begin{array}{c} \mathbf {P} _ {0} ^ {\prime} \xrightarrow {u _ {0}} \mathbf {P} _ {0} \\ p _ {0} ^ {\prime} \Big \downarrow \qquad \qquad \qquad \qquad \qquad \qquad \qquad \qquad \qquad \Big \downarrow p _ {0} \\ \mathbf {M} ^ {\prime} \xrightarrow [ f ]{} \mathbf {M}. \end{array}
\]

It then suffices to apply the induction hypothesis to the homomorphism

\[
\operatorname{Ker} p _ {0} ^ {\prime} \rightarrow \operatorname{Ker} p _ {0}
\]

deduced from \(u_{0}\) .

Lemma 5. --- Consider a commutative diagram

\[
\begin{array}{c} \mathbf {P} ^ {\prime} \xrightarrow {u} \mathbf {P} \\ p ^ {\prime} \Big \downarrow \quad \Big \downarrow p \\ 0 \longrightarrow \mathbf {M} ^ {\prime} \xrightarrow {f} \mathbf {M} \longrightarrow \mathbf {M} ^ {\prime \prime} \longrightarrow 0 \end{array}
\]

where (P, p) (resp. (P′, p′)) is a left resolution of M (resp. M′), and where the bottom horizontal row is an exact sequence. There exists a homologism \(p''\) : Con \((u) \to \mathbf{M}''\) .

Indeed, the exact sequence (X, p.~37, prop. 7)

\[
0 \to \mathrm{P} \xrightarrow {\pi} \operatorname{Con} (u) \xrightarrow {\delta} \mathrm{P} ^ {\prime} (- 1) \to 0
\]

gives an exact homology sequence

\[
\begin{array}{c}\rightarrow \mathrm{H} _ {n} (\mathrm{P}) \rightarrow \mathrm{H} _ {n} (\mathrm{Con} (u)) \rightarrow \mathrm{H} _ {n - 1} (\mathrm{P} ^ {\prime}) \rightarrow \dots\\\dots \rightarrow \mathrm{H} _ {1} (\mathrm{Con} (u)) \rightarrow \mathrm{H} _ {0} (\mathrm{P} ^ {\prime}) \stackrel {{\partial}} {{\rightarrow}} \mathrm{H} _ {0} (\mathrm{P}) \rightarrow \mathrm{H} _ {0} (\mathrm{Con} (u)) \rightarrow 0.\end{array}
\]

According to X, p.~38, lemma 3 a), we have \(\partial = -\mathrm{H}_{0}(u)\) . Since \(\mathrm{H}_{n}(\mathbf{P}) = 0 = \mathrm{H}_{n}(\mathbf{P}')\) for \(n > 0\) and \(\mathrm{H}_{0}(u): \mathrm{H}_{0}(\mathrm{P}') \to \mathrm{H}_{0}(\mathbf{P})\) is identified with \(f: \mathbf{M}' \to \mathbf{M}\) , we conclude that \(\mathrm{H}_{n}(\mathrm{Con}(u)) = 0\) for \(n > 0\) and that \(\mathrm{H}_{0}(\mathrm{Con}(u))\) is isomorphic to \(\mathbf{M}''\) , whence the lemma.

Now let us prove the theorem.

\begin{enumerate}
\def\labelenumi{\alph{enumi})}
\tightlist
\item
  Let M be an A-module of type C. For every left resolution \((\mathbf{P}, p)\) of M, bounded of type \(C_{0}\) , the element \(\chi_{\mathcal{C}_{0}}(\mathbf{P})\) of \(\mathbf{K}(\mathcal{C}_{0})\) depends only on M. Indeed, let \((\mathbf{P}_{1}, p_{1})\) and \((\mathbf{P}_{2}, p_{2})\) be two resolutions of this type. Consider the resolution
\end{enumerate}

\[
(\mathbf {P} _ {1} \times \mathbf {P} _ {2}, p _ {1} \times p _ {2})
\]

of the A-module \(\mathbf{M} \times \mathbf{M}\) and the homomorphism \(\Delta : x \mapsto (x, x)\) of \(\mathbf{M}\) to \(\mathbf{M} \times \mathbf{M}\) . By Lemma 4, there exists a resolution \((\mathbf{Q}, q)\) of \(\mathbf{M}\) bounded of type \(\mathcal{C}_0\) and a commutative diagram

\[
\begin{array}{c c c} \mathbf {Q} & \stackrel {{u}} {{\longrightarrow}} \mathbf {P} _ {1} & \times \mathbf {P} _ {2} \\ q \Big \downarrow & & \Big \downarrow p _ {1} \times p _ {2} \\ \mathbf {M} & \stackrel {{\Delta .}} {{\longrightarrow}} \mathbf {M} & \times \mathbf {M}; \end{array}
\]

we deduce a commutative diagram

\[
\begin{array}{c c c} \mathrm{Q} & \xrightarrow {u \circ p r _ {i}} & \mathrm{P} _ {i} \\ q \Big \downarrow & & \Big \downarrow p _ {i} \\ \mathbf {M} & \xrightarrow {\mathrm{l} _ {\mathrm{M}}} & \mathbf {M}, \end{array} \qquad i = 1, 2.
\]

By lemma 5, Con \((u \circ pr_i)\) has zero homology, hence \(u \circ pr_i\) is a homologism and \(\chi_{\mathcal{C}_0}(\mathbf{Q}) = \chi_{\mathcal{C}_0}(\mathbf{P}_i)\) (X, p.~41, prop. 10); it follows that \(\chi_{\mathcal{C}_0}(\mathbf{P}_1) = \chi_{\mathcal{C}_0}(\mathbf{P}_2)\) as announced.

\begin{enumerate}
\def\labelenumi{\alph{enumi})}
\setcounter{enumi}{1}
\tightlist
\item
  For every A-module M of type C, let \(\varphi(M) \in K(\mathcal{C}_{0})\) be the common value of \(\chi_{\mathcal{C}_{0}}(P)\) for all bounded left resolutions P of M of type \(C_{0}\) . Let us show that the function \(\varphi : \mathcal{C} \to K(\mathcal{C}_{0})\) is additive. Let therefore
\end{enumerate}

\[
0 \to \mathbf {M} ^ {\prime} \xrightarrow {f} \mathbf {M} \to \mathbf {M} ^ {\prime \prime} \to 0
\]

be an exact sequence of A-modules of type C. By Lemma 4, there exists a commutative diagram

\[
\begin{array}{c} \mathrm{P} ^ {\prime} \xrightarrow {u} \mathrm{P} \\ p ^ {\prime} \Big \downarrow \qquad \qquad \qquad \Big \downarrow p \\ 0 \longrightarrow \mathrm{M} ^ {\prime} \xrightarrow {f} \mathrm{M} \longrightarrow \mathrm{M} ^ {\prime \prime} \longrightarrow 0 \end{array}
\]

where (P, p) and (P', p') are bounded left resolutions of type \(\mathcal{C}_0\) . Then we have

\[
\varphi (\mathbf {M}) = \chi_ {\mathcal {C} _ {0}} (\mathbf {P}), \quad \varphi (\mathbf {M} ^ {\prime}) = \chi_ {\mathcal {C} _ {0}} (\mathbf {P} ^ {\prime})
\]

and by Lemma 5

\[
\varphi (\mathbf {M} ^ {\prime \prime}) = \chi_ {\mathcal {C} _ {0}} (\operatorname{Con} (u)) = \chi_ {\mathcal {C} _ {0}} (\mathrm{P}) - \chi_ {\mathcal {C} _ {0}} (\mathrm{P} ^ {\prime}) = \varphi (\mathrm{M}) - \varphi (\mathrm{M} ^ {\prime});
\]

which is what we wanted to prove.

\begin{enumerate}
\def\labelenumi{\alph{enumi})}
\setcounter{enumi}{2}
\tightlist
\item
  Let then \(\beta: \mathrm{K}(\mathcal{C}) \to \mathrm{K}(\mathcal{C}_0)\) be the homomorphism such that, with the preceding notation, we have \(\beta([\mathrm{M}]_{\mathcal{C}}) = \chi_{\mathcal{C}_0}(\mathrm{P})\) . Since \(p\) is a homologism, we have \(\chi_{\mathcal{C}}(\mathrm{P}) = [\mathrm{M}]_{\mathcal{C}}\) , hence \(\alpha \circ \beta([\mathrm{M}]_{\mathcal{C}}) = \alpha(\chi_{\mathcal{C}_0}(\mathrm{P})) = \chi_{\mathcal{C}}(\mathrm{P}) = [\mathrm{M}]_{\mathcal{C}}\) and \(\alpha \circ \beta = 1_{\mathrm{K}(\mathcal{C})}\) . If M is of type \(\mathcal{C}_0\) , then (M, \(1_{\mathrm{M}}\) ) is a resolution of M, hence \(\varphi(\mathrm{M}) = [\mathrm{M}]_{\mathcal{C}_0}\) and \(\beta \circ \alpha = 1_{\mathrm{K}(\mathcal{C}_0)}\) , which completes the proof.
\end{enumerate}

We shall apply this theorem to modules of "finite projective dimension" in § 8 (X, p.~137).

\section{TORSION PRODUCT}

In §§ 4 to 8, we denote by k a commutative ring, and by A an associative unital k-algebra. The role of k is auxiliary in nature; we have principally in view the following three particular cases:

\begin{enumerate}
\def\labelenumi{\alph{enumi})}
\item
  one considers an arbitrary ring A, sets k = Z and endows A with its natural structure of Z-algebra,
\item
  one considers an arbitrary ring A, takes for k the center of A,
\item
  one considers a commutative ring A and takes k = A.
\end{enumerate}

\subsection{Tensor product of two complexes}

Let (C, d) be a complex of right A-modules and (C', d') a complex of left A-modules.

Endow the k-module \(C \otimes_{A} C'\) with the grading such that

\[
(\mathrm{C} \otimes_ {\mathbf {A}} \mathrm{C} ^ {\prime}) _ {n} = \sum_ {p + q = n} (\mathrm{C} _ {p} \otimes \mathrm{C} _ {q} ^ {\prime})
\]

and let D denote the unique k-linear endomorphism of degree \((-1)\) of C \(\otimes_{A} C'\) such that

\[
\mathbf {D} (x \otimes x ^ {\prime}) = d x \otimes x ^ {\prime} + (- 1) ^ {p} x \otimes d ^ {\prime} x ^ {\prime}, \quad x \in \mathrm{C} _ {p}, y \in \mathrm{C} _ {q} ^ {\prime}, p, q \in \mathbf {Z}. \tag {1}
\]

One has \(D \circ D = 0\) since, with the notations of (1),

\[
\mathrm{D} ^ {2} (x \otimes x ^ {\prime}) = d d x \otimes x ^ {\prime} + (- 1) ^ {p - 1} d x \otimes d ^ {\prime} x ^ {\prime} + (- 1) ^ {p} d x \otimes d ^ {\prime} x ^ {\prime} - x \otimes d ^ {\prime} d ^ {\prime} x ^ {\prime}.
\]

The complex of \(k\) -modules \((\mathbf{C} \otimes_{\mathbf{A}} \mathbf{C}', \mathbf{D})\) is called the tensor product complex of the complexes \((\mathbf{C}, d)\) and \((\mathbf{C}', d')\) .

Remarks. --- 1) When \(\mathbf{C}'\) is reduced to \(\mathbf{C}_0' = \mathbf{M}\) , then \((\mathbf{C} \otimes_{\mathbf{A}} \mathbf{C}')_n = \mathbf{C}_n \otimes_{\mathbf{A}} \mathbf{M}\) and \(D = d \otimes 1_M\) ; for example \(\mathbf{C} \otimes_{\mathbf{A}} \mathbf{A}_s\) is naturally identified with \(C\) . Similarly, when \(C\) is reduced to \(\mathbf{C}_0 = P\) , then \((\mathbf{C} \otimes_{\mathbf{A}} \mathbf{C}')_n = P \otimes_{\mathbf{A}} \mathbf{C}_n'\) and \(D = 1_P \otimes d\) .

\begin{enumerate}
\def\labelenumi{\arabic{enumi})}
\setcounter{enumi}{1}
\tightlist
\item
  For every integer \(r\) , we have \((C \otimes_{A} C')(r) = C(r) \otimes_{A} C'\) , but \((C \otimes_{A} C')(r)\) and \(C \otimes_{A} C'(r)\) do not in general have the same differential.
\end{enumerate}

Let \(p, q\) be two integers, \(x \in Z_p(C)\) , \(x' \in Z_q(C')\) ; then the element \(x \otimes x'\) of \(C_p \otimes C_q'\) belongs to \(Z_{p+q}(C \otimes C')\) by (1); moreover, if \(y \in C_{p+1}\) , \(y' \in C_{q+1}'\) , we have

\[
(x + d y) \otimes (x ^ {\prime} + d ^ {\prime} y ^ {\prime}) = x \otimes x ^ {\prime} + \mathrm{D} (y \otimes x ^ {\prime} + (- 1) ^ {p} (x + d y) \otimes y ^ {\prime})
\]

by passing to quotients, we deduce a k-linear map, called canonical

\[
\gamma_ {p, q} (\mathrm{C}, \mathrm{C} ^ {\prime}): \mathrm{H} _ {p} (\mathrm{C}) \otimes_ {\mathrm{A}} \mathrm{H} _ {q} (\mathrm{C} ^ {\prime}) \rightarrow \mathrm{H} _ {p + q} (\mathrm{C} \otimes_ {\mathrm{A}} \mathrm{C} ^ {\prime});
\]

if we equip \(\mathbf{H}(\mathbf{C})\otimes_{\mathbf{A}}\mathbf{H}(\mathbf{C}^{\prime})\) with the grading such that

\[
\left(\mathrm{H} (\mathrm{C}) \otimes \mathrm{H} (\mathrm{C} ^ {\prime})\right) _ {n} = \sum_ {p + q = n} \mathrm{H} _ {p} (\mathrm{C}) \otimes_ {\mathrm{A}} \mathrm{H} _ {p} (\mathrm{C} ^ {\prime}),
\]

the \(\gamma_{p,q}\) define a graded \(k\) -linear map of degree 0

\[
\gamma (\mathbf {C}, \mathbf {C} ^ {\prime}): \mathrm{H} (\mathbf {C}) \otimes_ {\mathbf {A}} \mathrm{H} (\mathbf {C} ^ {\prime}) \rightarrow \mathrm{H} (\mathbf {C} \otimes_ {\mathbf {A}} \mathbf {C} ^ {\prime}).
\]

Proposition 6 of II, p.~59, is reformulated as follows:

PROPOSITION 1. --- If the complexes C and C' are zero on the right, then C ⊗\_A C' is zero on the right, and the canonical k-linear map

\[
\gamma_ {0, 0} (\mathbf {C}, \mathbf {C} ^ {\prime}): \mathrm{H} _ {0} (\mathbf {C}) \otimes_ {\mathbf {A}} \mathrm{H} _ {0} (\mathbf {C} ^ {\prime}) \rightarrow \mathrm{H} _ {0} (\mathbf {C} \otimes_ {\mathbf {A}} \mathbf {C} ^ {\prime})
\]

is bijective.

Let \(u: (\mathbf{C}, d) \to (\mathbf{C}_1, d_1)\) be a morphism of complexes of right A-modules and \(u': (\mathbf{C}', d') \to (\mathbf{C}_1', d_1')\) a morphism of complexes of left A-modules; then \(u \otimes u': \mathbf{C} \otimes_{\mathbf{A}} \mathbf{C}' \to \mathbf{C}_1 \otimes_{\mathbf{A}} \mathbf{C}_1'\) is a morphism of complexes of \(k\) -modules; indeed, it is graded of degree 0, and if one denotes by D and \(\mathbf{D}_1\) the differentials of \(\mathbf{C} \otimes \mathbf{C}'\) and \(\mathbf{C}_1 \otimes \mathbf{C}_1'\) we have, for \(p, q \in \mathbf{Z}, x \in \mathbf{C}_p, x' \in \mathbf{C}_q\) ,

\[
\begin{array}{c} (u \otimes u ^ {\prime}) (\mathrm{D} (x \otimes x ^ {\prime})) = u (d x) \otimes u ^ {\prime} (x ^ {\prime}) + (- 1) ^ {p} u (x) \otimes u ^ {\prime} (d ^ {\prime} x ^ {\prime}) = \\ = d _ {1} u (x) \otimes u ^ {\prime} (x ^ {\prime}) + (- 1) ^ {p} u (x) \otimes d _ {1} ^ {\prime} u ^ {\prime} (x ^ {\prime}) = \mathrm{D} _ {1} (u (x) \otimes u ^ {\prime} (x ^ {\prime}))  . \end{array}
\]

Moreover, the following diagram is commutative:

\[
\begin{array}{c} \mathrm{H(C)} \otimes_ {\mathrm{A}} \mathrm {H(C^ {\prime})} \xrightarrow {\gamma (\mathrm{C,C} ^ {\prime})} \mathrm{H(C} \otimes_ {\mathrm{A}} \mathrm {C^ {\prime}}) \\ \mathrm{H(u)} \otimes \mathrm {H(u^{\prime})} \Big \downarrow \\ \mathrm {H(C_ {1})} \otimes_ {\mathrm{A}} \mathrm {H(C_ {1} ^ {\prime})} \xrightarrow [ \gamma (\mathrm {C_ {1}} , \mathrm {C_ {1} ^ {\prime}}) ]{} \mathrm {H(C_ {1}} \otimes_ {\mathrm{A}} \mathrm {C_ {1} ^ {\prime}}). \end{array}
\]

Let A° be the opposite k-algebra of A, and let C° (resp. C'°) be the complex C (resp. C') viewed as a complex of left (resp. right) A°-modules.

\[
\sigma (\mathbf {C}, \mathbf {C} ^ {\prime}): \mathbf {C} \otimes_ {\mathbf {A}} \mathbf {C} ^ {\prime} \rightarrow \mathbf {C} ^ {\prime \circ} \otimes_ {\mathbf {A} ^ {\circ}} \mathbf {C} ^ {\circ}
\]

the unique graded \(k\) -linear map of degree 0 such that, for \(x \in C_p\) , \(x' \in C_q'\) , \(p, q \in \mathbf{Z}\) , one has

\[
\sigma (\mathbf {C}, \mathbf {C} ^ {\prime}) (x \otimes x ^ {\prime}) = (- 1) ^ {p q} x ^ {\prime} \otimes x.
\]

PROPOSITION 2. — The map \(\sigma(\mathbf{C},\mathbf{C}'):\mathbf{C}\otimes_{\mathbf{A}}\mathbf{C}'\to\mathbf{C}^{\prime\circ}\otimes_{\mathbf{A}^{\circ}}\mathbf{C}^{\circ}\) is an isomorphism of complexes of \(k\) -modules, whose inverse isomorphism is \(\sigma(\mathbf{C}^{\prime\circ},\mathbf{C}^{\circ})\) .

Since the maps \(\sigma(\mathbf{C},\mathbf{C}^{\prime})\) and \(\sigma(\mathbf{C}^{\circ},\mathbf{C}^{\circ})\) are inverses of each other, it suffices to prove that \(\sigma(\mathbf{C},\mathbf{C}^{\prime})\) is a morphism of complexes. Now, for

\[
x \in \mathrm{C} _ {p} = \mathrm{C} _ {p} ^ {\circ}, \quad x ^ {\prime} \in \mathrm{C} _ {q} ^ {\prime} = \mathrm{C} _ {q} ^ {\prime \circ}, \quad p, q \in \mathbf {Z},
\]

we have, writing \(D^{c}\) for the differential of \(C' \otimes_{A^{0}} C\) ,

\[
\begin{array}{l} \sigma (\mathrm{C}, \mathrm{C} ^ {\prime}) \circ \mathrm{D} (x \otimes x ^ {\prime}) = \sigma (\mathrm{C}, \mathrm{C} ^ {\prime}) (d x \otimes x ^ {\prime} + (- 1) ^ {p} x \otimes d ^ {\prime} x ^ {\prime}) = \\ = (- 1) ^ {(p + 1) q} x ^ {\prime} \otimes d x + (- 1) ^ {p + p (q + 1)} d ^ {\prime} x ^ {\prime} \otimes x = (- 1) ^ {p q} d ^ {\prime} x ^ {\prime} \otimes x + (- 1) ^ {p q + q} x ^ {\prime} \otimes d x \\ = (- 1) ^ {p q} \mathrm{D} ^ {\circ} (x ^ {\prime} \otimes x) = \mathrm{D} ^ {\circ} \circ \sigma (\mathrm{C}, \mathrm{C} ^ {\prime}) (x \otimes x ^ {\prime}); \end{array}
\]

this gives \(\sigma(C, C') \circ D = D^{\circ} \circ \sigma(C, C')\) , whence the desired assertion.

The isomorphism \(\sigma(\mathbf{C},\mathbf{C}^{\prime}):\mathbf{C}\otimes_{\mathbf{A}}\mathbf{C}^{\prime}\to\mathbf{C}^{\circ}\otimes_{\mathbf{A}^{\circ}}\mathbf{C}^{-}\) is called the commutation isomorphism of the tensor product of the complexes \(\mathbf{C}\) and \(\mathbf{C}^{\prime}\) .

If \(u: \mathbf{C} \to \mathbf{C}_{1}\) and \(v: \mathbf{C}' \to \mathbf{C}_{1}'\) are two morphisms of complexes as above, we have a commutative diagram:

\[
\begin{array}{c} \mathrm{C} \otimes_ {\mathrm{A}} \mathrm{C} ^ {\prime} \xrightarrow {\sigma (\mathrm{C} , \mathrm{C} ^ {\prime})} \mathrm{C} ^ {\prime \circ} \otimes_ {\mathrm{A} ^ {\circ}} \mathrm{C} ^ {\circ} \\ u \otimes u ^ {\prime} \Bigg \downarrow \\ \mathrm{C} _ {1} \otimes_ {\mathrm{A}} \mathrm{C} _ {1} ^ {\prime} \xrightarrow [ \sigma (\mathrm{C} _ {1} , \mathrm{C} _ {1} ^ {\prime}) ]{} \mathrm{C} _ {1} ^ {\prime \circ} \otimes_ {\mathrm{A} ^ {\circ}} \mathrm{C} _ {1} ^ {\circ}. \end{array}
\]

Suppose for the remainder of this number that the ring A is commutative (cf.~no. 9 for the general case).

Let C, C′, C′′ be three complexes of A-modules; the canonical homomorphism of A-modules (III, p. 64)

\[
\varphi : \left(\mathrm{C} \otimes_ {\mathrm{A}} \mathrm{C} ^ {\prime}\right) \otimes_ {\mathrm{A}} \mathrm{C} ^ {\prime \prime} \rightarrow \mathrm{C} \otimes_ {\mathrm{A}} \left(\mathrm{C} ^ {\prime} \otimes_ {\mathrm{A}} \mathrm{C} ^ {\prime \prime}\right)
\]

is an isomorphism of complexes, as one verifies immediately using the definitions.

More generally, let \((\mathbf{C}^{(i)}, d^{(i)})_{i \in I}\) be a family of complexes of A-modules, where the set I is finite and totally ordered; for simplicity of notation, we will identify I with the interval \([1, r]\) of N. Endow the A-module \(C = \bigotimes_{i=1}^{r} C^{(i)}\) with the grading such that

\[
\mathrm{C} _ {n} = \sum_ {p _ {1} + p _ {2} + \dots + p _ {r} = n} (\mathrm{C} ^ {(1)}) _ {p _ {1}} \otimes (\mathrm{C} ^ {(2)}) _ {p _ {2}} \otimes \dots \otimes (\mathrm{C} ^ {(r)}) _ {p _ {r}},
\]

and let us define a graded A-endomorphism of degree (-1) of C by

\[
\mathrm{D} \left(x _ {1} \otimes \dots \otimes x _ {r}\right) = \sum_ {j = 1} ^ {r} (- 1) ^ {p _ {1} + \dots + p _ {j - 1}} x _ {1} \otimes \dots \otimes x _ {j - 1} \otimes d _ {j} x _ {j} \otimes x _ {j + 1} \otimes \dots \otimes x _ {r}
\]

where \(x_{i}\in (\mathbf{C}^{(i)})_{p_{i}}\) for \(i = 1,\dots,n\) . Then (C, D) is a complex of A-modules called the tensor product complex of the family \((\mathbf{C}_{i}, d_{i})\) . For every strictly increasing sequence \(r_{0}, \ldots, r_{k}\) of [0, r] such that \(r_{0} = 0\) , \(r_{k} = r\) , the canonical associativity isomorphism

\[
\bigotimes_ {j = 0} ^ {k - 1} \left(\bigotimes_ {i = r _ {j} + 1} ^ {r _ {j + 1}} \mathbf {C} ^ {(i)}\right)\rightarrow \bigotimes_ {i = 1} ^ {r} \mathbf {C} ^ {(i)}
\]

is an isomorphism of complexes.

We define as above a graded homomorphism of degree 0

\[
\gamma ((\mathbf {C} ^ {(i)})): \underset {i \in \mathbf {I}} {\otimes} \mathrm{H} (\mathbf {C} ^ {(i)}) \to \mathrm{H} \left(\underset {i \in \mathbf {I}} {\otimes} \mathbf {C} ^ {(i)}\right) .
\]

Remarks. --- 3) One can define the tensor product of a finite family of complexes without endowing the index set with a total order (X, p.~185, exercise 3).

\begin{enumerate}
\def\labelenumi{\arabic{enumi})}
\setcounter{enumi}{3}
\tightlist
\item
  Suppose that each \(\mathbf{C}^{(i)}\) is endowed with a graded algebra structure compatible with its grading and such that the \(d^{(i)}\) are antiderivations (III, p.~117). Let us then equip \(\bigotimes_{i\in I}\mathbf{C}^{(i)}\) with the structure of left graded tensor product algebra of the given structures (III, p.~49). Then D is an antiderivation. Indeed, using the associativity of the tensor product, one may assume that I = \{1, 2\}; let then \(p_1\) , \(q_1\) , \(p_2\) , \(q_2 \in \mathbf{Z}\) , \(x_1 \in (\mathbf{C}^{(1)})_{p_1}\) , \(y_1 \in (\mathbf{C}^{(1)})_{q_1}\) , \(x_2 \in (\mathbf{C}^{(2)})_{p_2}\) , \(y_2 \in (\mathbf{C}^{(2)})_{q_2}\) ; we have
\end{enumerate}

\[
\begin{array}{l} \left(\mathrm{D} (x _ {1} \otimes x _ {2})\right) (y _ {1} \otimes y _ {2}) + (- 1) ^ {p _ {1} + p _ {2}} (x _ {1} \otimes x _ {2}) (\mathrm{D} (y _ {1} \otimes y _ {2})) = \\ = (d x _ {1} \otimes x _ {2} + (- 1) ^ {p _ {1}} x _ {1} \otimes d x _ {2}) (y _ {1} \otimes y _ {2}) + \\ \quad + (- 1) ^ {p _ {1} + p _ {2}} (x _ {1} \otimes x _ {2}) (d y _ {1} \otimes y _ {2} + (- 1) ^ {q _ {1}} y _ {1} \otimes d y _ {2}) = \\ = (- 1) ^ {p _ {2} q _ {1}} (d x _ {1}) y _ {1} \otimes x _ {2} y _ {2} + (- 1) ^ {p _ {1} + (p _ {2} - 1) q _ {1}} x _ {1} y _ {1} \otimes (d x _ {2}) y _ {2} + \\ \quad + (- 1) ^ {p _ {1} + p _ {2} + p _ {2} (q _ {1} - 1)} x _ {1} d y _ {1} \otimes x _ {2} y _ {2} + (- 1) ^ {p _ {1} + p _ {2} + q _ {1} + p _ {2} q _ {1}} x _ {1} y _ {1} \otimes x _ {2} d y _ {2} \\ = (- 1) ^ {p _ {2} q _ {1}} [ (d x _ {1}) y _ {1} + (- 1) ^ {p _ {1}} x _ {1} d y _ {1} ] \otimes x _ {2} y _ {2} + \\ \quad + (- 1) ^ {p _ {1} + q _ {1} + p _ {2} q _ {1}} x _ {1} y _ {1} \otimes ((d x _ {2}) y _ {2} + (- 1) ^ {p _ {2}} x _ {2} d y _ {2}) \\ = (- 1) ^ {p _ {2} q _ {1}} [ d (x _ {1} y _ {1}) \otimes x _ {2} y _ {2} + (- 1) ^ {p _ {1} + q _ {1}} x _ {1} y _ {1} \otimes d (x _ {2} y _ {2}) ] \\ = (- 1) ^ {p _ {2} q _ {1}} \mathrm{D} (x _ {1} y _ {1} \otimes x _ {2} y _ {2}) = \mathrm{D} ((x _ {1} \otimes x _ {2}) (y _ {1} \otimes y _ {2}))  . \end{array}
\]

\subsection{Tensor products and homotopy}

PROPOSITION 3. --- Let C, C \(_{1}\) be two complexes of right A-modules, C', C \(_{1}'\) two complexes of left A-modules, and u : C → C \(_{1}\) , v : C → C \(_{1}\) , u' : C' → C \(_{1}'\) , v' : C' → C \(_{1}'\) be morphisms of complexes.

\begin{enumerate}
\def\labelenumi{\alph{enumi})}
\item
  If \(u\) and \(u'\) are homotopic to \(v\) and \(v'\) respectively, then the two morphisms \(u \otimes u'\) and \(v \otimes v'\) from \(\mathbf{C} \otimes_{\mathbf{A}} \mathbf{C}'\) to \(\mathbf{C}_1 \otimes_{\mathbf{A}} \mathbf{C}_1'\) are homotopic.
\item
  If u and \(u'\) are homotopisms, \(u \otimes u'\) is a homotopism.
\item
  If C or \(\mathbf{C}'\) is homotopic to zero, \(\mathbf{C} \otimes_{\mathbf{A}} \mathbf{C}'\) is homotopic to zero.
\end{enumerate}

Let us denote by the same letter d the differentials of the complexes C, \(C_{1}\) , \(C'\) , \(C_{1}'\) and by D the differentials of the complexes \(C \otimes_{A} C'\) and \(C_{1} \otimes_{A} C_{1}'\) .

If \(u\) (resp. \(u'\) ) is homotopic to \(v\) (resp. \(v'\) ), there exists a graded homomorphism of degree 1 \(s: \mathbf{C} \to \mathbf{C}_1\) (resp. \(s': \mathbf{C}' \to \mathbf{C}_1'\) ) such that

\[
u - v = d s + s d \qquad (\mathrm{resp.} u ^ {\prime} - v ^ {\prime} = d s ^ {\prime} + s ^ {\prime} d).\tag{2}
\]

Let \(\mathbf{S}:\mathbf{C}\otimes_{\mathbf{A}}\mathbf{C}^{\prime}\to \mathbf{C}_{1}\otimes_{\mathbf{A}}\mathbf{C}_{1}^{\prime}\) be the unique graded homomorphism of degree 1 such that, for \(x\in \mathbf{C}_p\) , \(y\in \mathbf{C}_q^\prime\) , \(p,q\in \mathbf{Z}\) , one has

\[
\mathbf {S} (x \otimes y) = s (x) \otimes u ^ {\prime} (y) + (- 1) ^ {p} v (x) \otimes s ^ {\prime} (y).\tag{3}
\]

We then have, with the preceding notation:

\[
\begin{array}{r l} (\mathrm{DS} + \mathrm{SD}) (x \otimes y) & = \mathrm{D} (s x \otimes u ^ {\prime} y) + (- 1) ^ {p} \mathrm{D} (v x \otimes s ^ {\prime} y) + \mathrm{S} (d x \otimes y) + \\ & + (- 1) ^ {p} \mathrm{S} (x \otimes d y) = \\ & = d s x \otimes u ^ {\prime} y + (- 1) ^ {p + 1} s x \otimes d u ^ {\prime} y + (- 1) ^ {p} d v x \otimes s ^ {\prime} y + v x \otimes d s ^ {\prime} y + \\ & + s d x \otimes u ^ {\prime} y + (- 1) ^ {p - 1} v d x \otimes s ^ {\prime} y + (- 1) ^ {p} s x \otimes u ^ {\prime} d y + v x \otimes s ^ {\prime} d y \\ & = (d s + s d) (x) \otimes u ^ {\prime} y + v x \otimes (d s ^ {\prime} + s ^ {\prime} d) (y) \\ & = (u x - v x) \otimes u ^ {\prime} y + v x \otimes (u ^ {\prime} y - v ^ {\prime} y) = u x \otimes u ^ {\prime} y - v x \otimes v ^ {\prime} y. \end{array}
\]

This gives DS + SD = u ⊗ u' - v ⊗ v', whence a).

Let us prove b). If u and \(u'\) are homotopisms, there exist homomorphisms of complexes \(\alpha: C_{1} \to C\) and \(\alpha': C_{1}' \to C'\) such that \(u \circ \alpha\) , \(\alpha \circ u\) , \(u' \circ \alpha'\) , \(\alpha' \circ u'\) are homotopic respectively to \(Id_{C_{1}}\) , \(Id_{C}\) , \(Id_{C_{1}}\) and \(Id_{C'}\) . Then \((u \otimes u') \circ (\alpha \otimes \alpha')\) , which is equal to \((u \circ \alpha) \otimes (u' \circ \alpha')\) , is homotopic by a) to \(Id_{C_{1}} \otimes Id_{C_{1}} = Id_{C_{1} \otimes C_{1}}\) , whereas \((\alpha \otimes \alpha') \circ (u \otimes u')\) is homotopic to \(Id_{C \otimes C'}\) , whence b). Finally, c) follows from b) applied to the case where \(C_{1}\) or \(C_{1}'\) is zero.

COROLLARY 1. --- Let C' be a split complex of left A-modules such that H(C') is flat. For every complex C of right A-modules, the canonical map

\[
\gamma (\mathrm{C}, \mathrm{C} ^ {\prime}): \mathrm{H} (\mathrm{C}) \otimes_ {\mathrm{A}} \mathrm{H} (\mathrm{C} ^ {\prime}) \rightarrow \mathrm{H} (\mathrm{C} \otimes_ {\mathrm{A}} \mathrm{C} ^ {\prime})
\]

is bijective.

According to X, p.~35, def. 6, there exists a homotopism \(u': C' \to H(C')\) . By prop. 3, \(1_{C} \otimes u': C \otimes_{A} C' \to C \otimes_{A} H(C')\) is a homotopism; since

\[
\mathrm{H} \left(1 _ {\mathrm{C}} \otimes u ^ {\prime}\right) \circ \gamma (\mathrm{C}, \mathrm{C} ^ {\prime}) = \gamma (\mathrm{C}, \mathrm{H} (\mathrm{C} ^ {\prime})) \circ \left(1 _ {\mathrm{C}} \otimes \mathrm{H} (u ^ {\prime})\right),
\]

and \(\mathbf{H}(1_{\mathbf{C}} \otimes u')\) and \(\mathbf{H}(u')\) are bijective, it suffices to prove that \(\gamma(\mathbf{C}, \mathbf{H}(\mathbf{C}'))\) is bijective, and we are reduced to the case where \(\mathbf{C}'\) is flat and has zero differential. In this case the canonical exact sequences

\begin{enumerate}
\def\labelenumi{(\Roman{enumi})}
\tightlist
\item
\end{enumerate}

\[
0 \to \mathbf {Z} (\mathbf {C}) \xrightarrow {i} \mathbf {C} \xrightarrow {\delta} \mathbf {B} (\mathbf {C}) \to 0\tag{II}
\]

\[
0 \rightarrow \mathrm{B} (\mathrm{C}) \xrightarrow {j} \mathrm{Z} (\mathrm{C}) \xrightarrow {\pi} \mathrm{H} (\mathrm{C}) \rightarrow 0
\]

give exact sequences:

\[
\begin{array}{l} 0 \longrightarrow Z (C) \otimes_ {A} C ^ {\prime} \xrightarrow {i \otimes 1} C \otimes_ {A} C ^ {\prime} \quad \xrightarrow {\delta \otimes 1} B (C) \otimes_ {A} C ^ {\prime} \longrightarrow 0 \\ 0 \longrightarrow B (C) \otimes_ {A} C ^ {\prime} \xrightarrow {j \otimes 1} Z (C) \otimes_ {A} C ^ {\prime} \xrightarrow {\pi \otimes 1} H (C) \otimes_ {A} C ^ {\prime} \longrightarrow 0. \end{array}
\]

Since \(d = i \circ j \circ \delta\) , we have \(\mathbf{D} = d \otimes 1_{\mathbf{C}'} = (i \otimes 1) \circ (j \otimes 1) \circ (\delta \otimes 1)\), which shows that the canonical maps \(\mathbf{Z}(\mathbf{C}) \otimes_{\mathbf{A}} \mathbf{C}' \to \mathbf{Z}(\mathbf{C} \otimes_{\mathbf{A}} \mathbf{C}')\) and \(\mathbf{B}(\mathbf{C}) \otimes_{\mathbf{A}} \mathbf{C}' \to \mathbf{B}(\mathbf{C} \otimes_{\mathbf{A}} \mathbf{C}')\) are bijective, hence also \(\gamma(\mathbf{C}, \mathbf{C}')\) by passing to quotients.

COROLLARY 2. --- Let N be a flat left A-module. For every complex C of right A-modules, the canonical homomorphisms

\[
\gamma_ {n} (\mathbf {C}, \mathbf {N}): \mathrm{H} _ {n} (\mathbf {C}) \otimes_ {\mathbf {A}} \mathbf {N} \rightarrow \mathrm{H} _ {n} (\mathbf {C} \otimes_ {\mathbf {A}} \mathbf {N})
\]

are bijective.

COROLLARY 3. --- Let C' be a complex of left A-modules such that B(C') and H(C') are projective. For every complex C of right A-modules, the map γ(C, C') is bijective.

Indeed, C' is split (X, p.~35, example 4) and H(C') is projective; one can therefore apply Corollary 1.

Remarks. --- 1) Using the commutation isomorphisms, one deduces from Corollaries 1, 2, and 3 the analogous statements obtained by interchanging the roles of the two arguments of the tensor products.

\begin{enumerate}
\def\labelenumi{\arabic{enumi})}
\setcounter{enumi}{1}
\tightlist
\item
  We shall see below (X, p.~79, cor. 4) that the conclusion of cor. 1 is also true when one assumes that C' and H(C') are flat and C' is bounded on the right.
\end{enumerate}

\subsection{Tensor product by a flat complex bounded on the right}

Lemma 1. --- Let C be a complex of right A-modules and E a complex of left A-modules. Suppose that H(C) = 0 and that E is flat and bounded on the right. Then H(C ⊗\_A E) = 0.

For \(k \in \mathbf{Z}\) , let \(\mathbf{T}^{(k)}\) be the subcomplex of \(\mathbf{C} \otimes_{\mathbf{A}} \mathbf{E}\) such that

\[
\mathrm{T}_{n}^{(k)} = \sum_{\substack{p + q = n\\ q\leqslant k}}\mathrm{C}_{p}\otimes_{\mathrm{A}}\mathrm{E}_{q};
\]

then \(\mathbf{T}^{(k - 1)}\subset \mathbf{T}^{(k)}\) and we have an exact sequence of complexes

\[
0 \longrightarrow \mathrm{T} ^ {(k - 1)} \stackrel {i _ {k}} {\longrightarrow} \mathrm{T} ^ {(k)} \stackrel {\pi} {\longrightarrow} \mathrm{C} \otimes_ {\mathbf {A}} \mathrm{E} _ {k} (- k) \longrightarrow 0
\]

where \(i_{k}\) is the canonical injection and where \(\pi\) projects the preceding direct sum onto its factor \(\mathbf{C}_{n-k} \otimes_{\mathbf{A}} \mathbf{E}_{k} = (\mathbf{C} \otimes_{\mathbf{A}} \mathbf{E}_{k}(-k))_{n}\) . By cor. 2 above, we have \(\mathbf{H}(\mathbf{C} \otimes_{\mathbf{A}} \mathbf{E}_{k}(-k)) = 0\) , hence \(i_{k}\) is a homologism. We have \(\mathbf{T}^{(k)} = 0\) for \(k\) small enough, since E is bounded on the right, hence \(\mathrm{H}(\mathrm{T}^{(k)}) = 0\) for all k by induction on k. Finally, the canonical morphism \(\lim_{\mathbf{A}} \mathrm{T}^{(k)} \to \mathbf{C} \otimes_{\mathbf{A}} \mathbf{E}\) is obviously an isomorphism, hence \(\mathrm{H}(\mathbf{C} \otimes_{\mathbf{A}} \mathbf{E}) = 0\) (X, p.~28, prop. 1).

Lemma 2. --- If \(u: \mathbf{C} \to \mathbf{C}'\) is a morphism of complexes of right A-modules and E is a complex of left A-modules, then the complexes Con \((u) \otimes_{\mathbf{A}} \mathbf{E}\) and Con \((u \otimes 1_{\mathbf{E}})\) are isomorphic.

By definition, Con \((u) \otimes_{\mathbf{A}} \mathbf{E}\) is the graded module \((\mathbf{C}'(-1) \oplus \mathbf{C}) \otimes_{\mathbf{A}} \mathbf{E}\) endowed with the differential D such that, for \(x \in \mathbf{C}_p\) , \(y' \in \mathbf{C}'(-1)_p = \mathbf{C}_{p-1}'\) , \(z \in \mathbf{E}_q\) , one has

\[
\mathrm{D} ((y ^ {\prime}, x) \otimes z) = (- d y ^ {\prime}, d x - u (y ^ {\prime})) \otimes z + (- 1) ^ {p} (y ^ {\prime}, x) \otimes d z,\tag{4}
\]

while \(\operatorname{Con}(u \otimes 1_{\mathbf{E}})\) is the graded module \((\mathbf{C}' \otimes_{\mathbf{A}} \mathbf{E})(-1) \oplus (\mathbf{C} \otimes_{\mathbf{A}} \mathbf{E})\) equipped with the differential \(\mathbf{D}_1\) such that, for \(x \in \mathbf{C}_p\) , \(y' \in \mathbf{C}_{p-1}'\) , \(z \in \mathbf{E}_q\) , one has

\[
\mathrm{D} _ {1} \left(y ^ {\prime} \otimes z, x \otimes z\right) = (- d y ^ {\prime} \otimes z - (- 1) ^ {p - 1} y ^ {\prime} \otimes d z, d x \otimes z + (- 1) ^ {p} x \otimes d z - u \left(y ^ {\prime}\right) \otimes z)
\]

\[
= \left(- d y ^ {\prime} \otimes z, (d x - u (y ^ {\prime})) \otimes z\right) + (- 1) ^ {p} \left(y ^ {\prime} \otimes d z, x \otimes d z\right)
\]

whence the assertion.

PROPOSITION 4. — Let \(u: \mathbf{C} \to \mathbf{C}'\) be a homologism of complexes of right A-modules, and E a complex of left A-modules, flat and bounded on the right. Then

\[
u \otimes 1 _ {\mathbf {E}}: \mathbf {C} \otimes_ {\mathbf {A}} \mathbf {E} \rightarrow \mathbf {C} ^ {\prime} \otimes_ {\mathbf {A}} \mathbf {E}
\]

is a homologism of complexes of k-modules.

Indeed, according to X, p.~38, cor., \(u\) (resp. \(u \otimes 1_{\mathrm{E}}\) ) is a homologism if and only if \(\mathrm{H}(\mathrm{Con}(u)) = 0\) (resp. \(\mathrm{H}(\mathrm{Con}(u \otimes 1_{\mathrm{E}})) = 0\) ). We then conclude by Lemmas 1 and 2.

Remark. --- Using the commutation isomorphisms, one deduces from the preceding statements the analogous statements obtained by interchanging the roles of the two arguments of the tensor products.

\subsection{Definition and first properties of the torsion product}

For every A-module E, we denote by \(p_{E}: L(E) \to E\) the canonical free resolution of E (X, p.~50).

DEFINITION 1. --- Let M be a right A-module and N a left A-module. The torsion product of M and N is the graded k-module

\[
\operatorname{Tor} ^ {\mathbf {A}} (\mathbf {M}, \mathbf {N}) = \mathrm{H} (\mathrm{L} (\mathbf {M}) \otimes_ {\mathbf {A}} \mathrm{L} (\mathbf {N}))  .\tag{4}
\]

The homogeneous components of Tor \(^{A}\) (M, N) are denoted

\[
\operatorname{Tor} _ {n} ^ {\mathbf {A}} (\mathbf {M}, \mathbf {N}) = \mathrm{H} _ {n} (\mathrm{L} (\mathbf {M}) \otimes_ {\mathbf {A}} \mathrm{L} (\mathbf {N})).\tag{5}
\]

Since L(M) and L(N) are zero on the right, we have

\[
\operatorname{Tor} _ {n} ^ {\mathbf {A}} (\mathbf {M}, \mathbf {N}) = 0 \quad \text { for } \quad n <   0.\tag{6}
\]

Remark 1. --- We shall see below (X, p.~107, prop. 6) finiteness properties of the Tor \(^{A}\) (M, N) modules. For example, if A is commutative Noetherian and M and N are finitely generated A-modules, then each A-module Tor \(_{n}^{A}\) (M, N) is finitely generated.

Let \(f: \mathbf{M} \to \mathbf{M}'\) be a homomorphism of right A-modules and \(g: \mathbf{N} \to \mathbf{N}'\) be a homomorphism of left A-modules; we set \(\operatorname{Tor}^{\mathbf{A}}(f, g) = \operatorname{H}\left(\operatorname{L}(f) \otimes_{\mathbf{A}} \operatorname{L}(g)\right)\) ; it is a homomorphism of graded \(k\) -modules

\[
\operatorname{Tor} ^ {\mathbf {A}} (f, g): \operatorname{Tor} ^ {\mathbf {A}} (\mathbf {M}, \mathbf {N}) \rightarrow \operatorname{Tor} ^ {\mathbf {A}} (\mathbf {M} ^ {\prime}, \mathbf {N} ^ {\prime})
\]

whose homogeneous components are denoted

\[
\operatorname{Tor} _ {n} ^ {\mathbf {A}} (f, g): \operatorname{Tor} _ {n} ^ {\mathbf {A}} (\mathbf {M}, \mathbf {N}) \rightarrow \operatorname{Tor} _ {n} ^ {\mathbf {A}} (\mathbf {M} ^ {\prime}, \mathbf {N} ^ {\prime}).
\]

By prop. 1 of X, p.~62, the canonical homomorphism

\[
\gamma_ {0, 0}: \mathrm{H} _ {0} (\mathrm{L} (\mathrm{M})) \otimes_ {\mathrm{A}} \mathrm{H} _ {0} (\mathrm{L} (\mathrm{N})) \rightarrow \mathrm{H} _ {0} (\mathrm{L} (\mathrm{M}) \otimes_ {\mathrm{A}} \mathrm{L} (\mathrm{N}))
\]

is bijective; using the isomorphisms \(\mathbf{M} \to \mathbf{H}_{0}(\mathbf{L}(\mathbf{M}))\) and \(\mathbf{N} \to \mathbf{H}_{0}(\mathbf{L}(\mathbf{N}))\) we derive an isomorphism, called canonical

\[
\gamma_ {\mathbf {M}, \mathbf {N}}: \mathbf {M} \otimes_ {\mathbf {A}} \mathbf {N} \rightarrow \operatorname{Tor} _ {0} ^ {\mathbf {A}} (\mathbf {M}, \mathbf {N}).\tag{7}
\]

We shall always identify \(\operatorname{Tor}_{0}^{\mathbf{A}}(\mathbf{M},\mathbf{N})\) with \(\mathbf{M}\otimes_{\mathbf{A}}\mathbf{N}\) by this isomorphism. Then the \(k\) -linear map \(\operatorname{Tor}_{0}^{\mathbf{A}}(f,g)\) is identified with \(f\otimes g\) .

Remark 2. --- The morphism of complexes \(p_{\mathbf{M}} \otimes p_{\mathbf{N}}: \mathbf{L}(\mathbf{M}) \otimes_{\mathbf{A}} \mathbf{L}(\mathbf{N}) \to \mathbf{M} \otimes_{\mathbf{A}} \mathbf{N}\) induces on the homology of degree 0 the isomorphism

\[
\gamma_ {\mathbf {M}, \mathbf {N}} ^ {- 1}: \operatorname{Tor} _ {0} ^ {\mathbf {A}} (\mathbf {M}, \mathbf {N}) \rightarrow \mathbf {M} \otimes_ {\mathbf {A}} \mathbf {N}
\]

inverse of \(\gamma_{M,N}\) .

One has \(L(1_{\mathbf{M}}) = 1_{L(\mathbf{M})}\) , \(L(1_{\mathbf{N}}) = i_{L(\mathbf{N})}\) hence by passing to homology:

\[
\operatorname{Tor} ^ {\mathbf {A}} \left(1 _ {\mathbf {M}}, 1 _ {\mathbf {N}}\right) = 1 _ {\operatorname{Tor} ^ {\mathbf {A}} (\mathbf {M}, \mathbf {N})}.\tag{8}
\]

If \(f': \mathbf{M}' \to \mathbf{M}''\) (resp. \(g': \mathbf{N}' \to \mathbf{N}''\) ) is a homomorphism of right (resp. left) A-modules, we have \(\mathbf{L}(g' \circ g) = \mathbf{L}(g') \circ \mathbf{L}(g)\) and \(\mathbf{L}(f' \circ f) = \mathbf{L}(f') \circ \mathbf{L}(f)\) , hence

\[
\operatorname{Tor} ^ {\mathbf {A}} \left(f ^ {\prime} \circ f, g ^ {\prime} \circ g\right) = \operatorname{Tor} ^ {\mathbf {A}} \left(f ^ {\prime}, g ^ {\prime}\right) \circ \operatorname{Tor} ^ {\mathbf {A}} (f, g).\tag{9}
\]

Consider the morphisms of k-complexes

\[
\mathrm{L} (\mathrm{M}) \otimes_ {\mathrm{A}} \mathrm{N} \xleftarrow {1 \otimes p _ {\mathrm{N}}} \mathrm{L} (\mathrm{M}) \otimes_ {\mathrm{A}} \mathrm{L} (\mathrm{N}) \xrightarrow {p _ {\mathrm{M}} \otimes 1} \mathrm{M} \otimes_ {\mathrm{A}} \mathrm{L} (\mathrm{N})
\]

and the k-homomorphisms they induce in homology:

\[
\mathrm{H} (\mathrm{L} (\mathrm{M}) \otimes_ {\mathrm{A}} \mathrm{N}) \xleftarrow {\psi_ {\mathrm{M}} (\mathrm{N})} \operatorname{Tor} ^ {\mathrm{A}} (\mathrm{M}, \mathrm{N}) \xrightarrow {\overline {{\psi}} _ {\mathrm{N}} (\mathrm{M})} \mathrm{H} (\mathrm{M} \otimes_ {\mathrm{A}} \mathrm{L} (\mathrm{N}));
\]

by prop. 4 of X, p.~67, \(1 \otimes p_{\mathbf{N}}\) and \(p_{\mathbf{M}} \otimes 1\) are homologisms. Hence:

Proposition 5. --- The k-homomorphisms

\[
\begin{array}{l} \psi_ {\mathrm{M}} (\mathrm{N}): \operatorname{Tor} ^ {\mathrm{A}} (\mathrm{M}, \mathrm{N}) \to \mathrm{H} (\mathrm{L} (\mathrm{M}) \otimes_ {\mathrm{A}} \mathrm{N}) \\ \overline {{\psi}} _ {\mathrm{N}} (\mathrm{M}): \operatorname{Tor} ^ {\mathrm{A}} (\mathrm{M}, \mathrm{N}) \to \mathrm{H} (\mathrm{M} \otimes_ {\mathrm{A}} \mathrm{L} (\mathrm{N})) \end{array}
\]

are bijective.

COROLLARY. --- If M or N is flat, \(\operatorname{Tor}_{i}^{\mathbf{A}}(\mathbf{M},\mathbf{N}) = 0\) for \(i \geqslant 0\) .

Suppose N (resp. M) is flat; then \(p_{\mathbf{M}} \otimes 1 : \mathbf{L}(\mathbf{M}) \otimes_{\mathbf{A}} \mathbf{N} \to \mathbf{M} \otimes_{\mathbf{A}} \mathbf{N}\) (resp. \(1 \otimes p_{\mathbf{N}} : \mathbf{M} \otimes_{\mathbf{A}} \mathbf{L}(\mathbf{N}) \to \mathbf{M} \otimes_{\mathbf{A}} \mathbf{N}\) ) is a homologism (X, p.~67, prop. 4), hence \(\mathrm{H}_{i}(\mathrm{L}(\mathrm{M}) \otimes_{\mathbf{A}} \mathrm{N})\) (resp. \(\mathrm{H}_{i}(\mathrm{M} \otimes_{\mathbf{A}} \mathrm{L}(\mathrm{N}))\) ) is zero for i \textgreater{} 0.

Remark 3. --- If \(g: \mathbf{N} \to \mathbf{N}'\) is a homomorphism of left A-modules, then

\[
\left(1 _ {\mathrm{L} (\mathbf {M})} \otimes g\right) \circ \left(1 _ {\mathrm{L} (\mathbf {M})} \otimes 1 _ {\mathrm{N}}\right) = \left(1 _ {\mathrm{L} (\mathbf {M})} \otimes 1 _ {\mathrm{N}}\right) \circ \left(1 _ {\mathrm{L} (\mathbf {M})} \otimes \mathrm{L} (g)\right),
\]

therefore the diagram

\[
\begin{array}{c c c} \operatorname{Tor} ^ {\mathbf {A}} (\mathbf {M}, \mathbf {N}) & \xrightarrow {\psi_ {\mathbf {M}} (\mathbf {N})} & \mathrm{H} (\mathrm{L} (\mathbf {M}) \otimes_ {\mathbf {A}} \mathbf {N}) \\ \operatorname{Tor} ^ {\mathbf {A}} (1, g) \Big \downarrow & & \Big \downarrow \mathrm{H} (1 \otimes g) \\ \operatorname{Tor} ^ {\mathbf {A}} (\mathbf {M}, \mathbf {N} ^ {\prime}) & \xrightarrow {\psi_ {\mathbf {M}} (\mathbf {N} ^ {\prime})} & \mathrm{H} (\mathrm{L} (\mathbf {M}) \otimes_ {\mathbf {A}} \mathbf {N} ^ {\prime}) \end{array}
\]

is commutative.

Likewise, if \(f: \mathbf{M} \to \mathbf{M}'\) is a homomorphism of right A-modules, we have a commutative diagram:

\[
\begin{array}{c c} \operatorname{Tor} ^ {\mathbf {A}} (\mathrm{M}, \mathrm{N}) & \xrightarrow {\overline {{\psi}} _ {\mathrm{N}} (\mathrm{M})} \mathrm{H} (\mathrm{M} \otimes_ {\mathrm{A}} \mathrm{L} (\mathrm{N})) \\ \operatorname{Tor} ^ {\mathbf {A}} (f, 1) \Big \downarrow & \Big \downarrow \mathrm{H} (f \otimes 1) \\ \operatorname{Tor} ^ {\mathbf {A}} (\mathrm{M} ^ {\prime}, \mathrm{N}) & \xrightarrow {\overline {{\psi}} _ {\mathrm{N}} (\mathrm{M} ^ {\prime})} \mathrm{H} (\mathrm{M} ^ {\prime} \otimes_ {\mathrm{A}} \mathrm{L} (\mathrm{N})). \end{array}
\]

PROPOSITION 6. --- The map (f, g) ↦ Tor \(^{A}\) (f, g) :

\[
\operatorname{Hom} _ {\mathrm{A}} \left(\mathrm{M}, \mathrm{M} ^ {\prime}\right) \times \operatorname{Hom} _ {\mathrm{A}} \left(\mathrm{N}, \mathrm{N} ^ {\prime}\right)\rightarrow \operatorname{Hom} _ {k} \left(\operatorname{Tor} ^ {\mathrm{A}} (\mathrm{M}, \mathrm{N}), \operatorname{Tor} ^ {\mathrm{A}} \left(\mathrm{M} ^ {\prime}, \mathrm{N} ^ {\prime}\right)\right)
\]

is k-bilinear.

Let \(f \in \operatorname{Hom}_{\mathbf{A}}(\mathbf{M}, \mathbf{M}')\) , \(g_1, g_2 \in \operatorname{Hom}_{\mathbf{A}}(\mathbf{N}, \mathbf{N}')\) , \(\lambda_1, \lambda_2 \in k\) . Then the morphisms

\[
\lambda_ {1} (\mathrm{L} (f) \otimes g _ {1}) + \lambda_ {2} (\mathrm{L} (f) \otimes g _ {2}) \quad \text { and } \quad \mathrm{L} (f) \otimes (\lambda_ {1} g _ {1} + \lambda_ {2} g _ {2})
\]

of L(M) ⊗\_A N in L(M) ⊗\_A N′ coincide; by prop. 5 and remark 3, we therefore have

\[
\operatorname{Tor} ^ {\mathbf {A}} \left(f, \lambda_ {1} g _ {1} + \lambda_ {2} g _ {2}\right) = \lambda_ {1} \operatorname{Tor} ^ {\mathbf {A}} \left(f, g _ {1}\right) + \lambda_ {2} \operatorname{Tor} ^ {\mathbf {A}} \left(f, g _ {2}\right).\tag{10}
\]

We reason similarly for the mapping \(f\mapsto\mathrm{Tor}^{\mathbf{A}}(f,g)\) .

COROLLARY. --- Let \(\lambda\in k\) . If \(\lambda\) annihilates M or N, it annihilates Tor \(^{A}\) (M, N).

Indeed, \(\lambda .1_{\mathrm{Tor}(\mathbf{M},\mathbf{N})} = \mathrm{Tor}(\lambda .1_{\mathbf{M}},1_{\mathbf{N}}) = \mathrm{Tor}(1_{\mathbf{M}},\lambda .1_{\mathbf{N}})\) .

PROPOSITION 7. --- Let I and J be two sets, \((\mathbf{M}_{\alpha})_{\alpha \in I}\) a family of right A-modules, \((\mathbf{N}_{\beta})_{\beta \in J}\) a family of left A-modules. The homomorphism

\[
\bigoplus_ {\alpha \in I, \beta \in J} \operatorname{Tor} ^ {A} \left(M _ {\alpha}, N _ {\beta}\right)\rightarrow \operatorname{Tor} ^ {A} \left(\bigoplus_ {\alpha \in I} M _ {\alpha}, \bigoplus_ {\beta \in J} N _ {\beta}\right)
\]

deduced from the canonical homomorphisms \(M_{\alpha} \rightarrow \oplus M_{\alpha}\) and \(N_{\beta} \rightarrow \oplus N_{\beta}\) is bijective. It suffices to prove that for every right module M (resp. every left module N), the canonical homomorphism

\[
\bigoplus_ {\beta \in J} \operatorname{Tor} ^ {A} (M, N _ {\beta}) \rightarrow \operatorname{Tor} ^ {A} (M, \bigoplus_ {\beta \in J} N _ {\beta})
\]

\((\text{resp. } \bigoplus_{\alpha \in I} \text{Tor}^{\mathbf{A}}(\mathbf{M}_{\alpha}, \mathbf{N}) \to \text{Tor}^{\mathbf{A}} (\bigoplus_{\alpha \in I} \mathbf{M}_{\alpha}, \mathbf{N}))\) is bijective. This follows from what precedes, from Proposition 1 of X, p.~28, and from the canonical isomorphisms:

\[
\begin{array}{l} \underset {\beta} {\oplus} (\mathrm{L(M)} \otimes_ {\mathrm{A}} \mathrm {N}_ {\beta}) \to \mathrm{L(M)} \otimes_ {\mathrm{A}} (\underset {\beta} {\oplus} \mathrm {N}_ {\beta}), \\ \underset {\alpha} {\oplus} (\mathrm {M}_ {\alpha} \otimes_ {\mathrm{A}} \mathrm{L(N)}) \to (\underset {\alpha} {\oplus} \mathrm {M}_ {\alpha}) \otimes_ {\mathrm{A}} \mathrm{L(N)}. \end{array}
\]

An analogous argument gives:

PROPOSITION 8. --- Let I (resp. J) be a preordered set filtered to the right, \(((M_{\alpha}), (u_{\alpha' \alpha}))\) (resp. \(((N_{\beta}), (v_{\beta' \beta}))\) ) an inductive system of right (resp. left) A-modules relative to I (resp. J). The homomorphism of graded k-modules

\[
\lim _ {(\alpha , \beta) \in I \times J} \operatorname{Tor} ^ {A} \left(M _ {\alpha}, N _ {\beta}\right)\rightarrow \operatorname{Tor} ^ {A} \left(\lim _ {\alpha \in I} M _ {\alpha}, \lim _ {\beta \in J} N _ {\beta}\right),
\]

deduced from canonical A-homomorphisms \(\mathbf{M}_{\alpha} \to \varinjlim \mathbf{M}_{\alpha}\) and \(\mathbf{N}_{\beta} \to \varinjlim \mathbf{N}_{\beta}\) is bijective.

In particular, taking \(J = I\) and noting that the \((\alpha, \alpha), \alpha \in I\) form a cofinal subset of \(I \times I\) , we obtain:

COROLLARY. --- Let I be a preordered set filtered to the right, \((\mathbf{M}_{i}, u_{ji})\) (resp. \((\mathbf{N}_{i}, v_{ji})\) ) an inductive system of right (resp. left) A-modules relative to I. The homomorphism of graded k-modules

\[
\varinjlim_ {i \in I} \operatorname{Tor} ^ {A} (M _ {i}, N _ {i}) \rightarrow \operatorname{Tor} ^ {A} (\varinjlim_ {i \in I} M _ {i}, \varinjlim_ {i \in I} N _ {i}),
\]

deduced from canonical A-homomorphisms \(\mathbf{M}_i\to \varinjlim \mathbf{M}_i\) and \(\mathbf{N}_j\to \varinjlim \mathbf{N}_j\) is bijective.

Let M be a right A-module, N a left A-module, A° the opposite ring of A, M° the left A°-module underlying M, N° the right A°-module underlying M. We have L(M°) = L(M)° and L(N°) = L(N)°, whence a commutation isomorphism (X, p.~63, prop. 2)

\[
\sigma (\mathrm{L} (\mathrm{M}), \mathrm{L} (\mathrm{N})): \mathrm{L} (\mathrm{M}) \otimes_ {\mathrm{A}} \mathrm{L} (\mathrm{N}) \rightarrow \mathrm{L} (\mathrm{N} ^ {\circ}) \otimes_ {\mathrm{A} ^ {0}} \mathrm{L} (\mathrm{M} ^ {\circ}).
\]

Passing to homology, \(\sigma(\mathrm{L}(\mathbf{M}),\mathrm{L}(\mathbf{N}))\) induces a graded isomorphism of degree 0 \(\sigma_{\mathbf{M},\mathbf{N}}:\operatorname{Tor}^{\mathbf{A}}(\mathbf{M},\mathbf{N})\to\operatorname{Tor}^{\mathbf{A}^{\circ}}(\mathbf{N}^{\circ},\mathbf{M}^{\circ})\) called the commutation isomorphism of torsion products.

Note that \(\sigma_{\mathbf{N}^{\circ},\mathbf{M}^{\circ}}\circ\sigma_{\mathbf{M},\mathbf{N}}=\mathrm{Id}_{\mathrm{Tor}\left(\mathbf{M},\mathbf{N}\right)}\) and that \(\sigma_{\mathbf{M},\mathbf{N}}\) induces on the degree 0 terms the commutation homomorphism of the tensor product. On the other hand, if \(f:\mathbf{M}\to\mathbf{M}'\) and \(g:\mathbf{N}\to\mathbf{N}'\) are homomorphisms of A-modules, we have

\[
\operatorname{Tor} ^ {\mathbf {A} ^ {\circ}} (g, f) \circ \sigma_ {\mathbf {M}, \mathbf {N}} = \sigma_ {\mathbf {M} ^ {\prime}, \mathbf {N} ^ {\prime}} \circ \operatorname{Tor} ^ {\mathbf {A}} (f, g).
\]

\subsection{Connecting homomorphisms and exact sequences}

Let M be a right A-module. Recall that for every left A-module N, we defined in the preceding number (X, p.~69, prop. 5) an isomorphism

Let

\[
\psi_ {\mathbf {M}} (\mathbf {N}): \operatorname{Tor} ^ {\mathbf {A}} (\mathbf {M}, \mathbf {N}) \rightarrow \mathrm{H} (\mathrm{L} (\mathbf {M}) \otimes_ {\mathbf {A}} \mathbf {N}).\tag{ε}
\]

\[
0 \rightarrow \mathbf {N} ^ {\prime} \xrightarrow {u} \mathbf {N} \xrightarrow {v} \mathbf {N} ^ {\prime \prime} \rightarrow 0
\]

an exact sequence of left A-modules; the sequence of \(k\) -complexes

\[
\left(^ {\mathrm{M}} \mathcal {E}\right) \quad 0 \longrightarrow \mathrm{L} (\mathrm{M}) \otimes_ {\mathrm{A}} \mathrm{N} ^ {\prime} \xrightarrow {1 \otimes u} \mathrm{L} (\mathrm{M}) \otimes_ {\mathrm{A}} \mathrm{N} \xrightarrow {1 \otimes v} \mathrm{L} (\mathrm{M}) \otimes_ {\mathrm{A}} \mathrm{N} ^ {\prime \prime} \longrightarrow 0
\]

is then exact (X, p.~66, lemma 1); let

\[
\hat {c} \left(^ {\mathrm{M}} \mathcal {E}\right): \mathrm{H} (\mathrm{L} (\mathrm{M}) \otimes_ {\mathrm{A}} \mathrm{N} ^ {\prime \prime}) \rightarrow \mathrm{H} (\mathrm{L} (\mathrm{M}) \otimes_ {\mathrm{A}} \mathrm{N} ^ {\prime})
\]

the corresponding connecting homomorphism (X, p.~29).

DEFINITION 2. --- A connecting homomorphism of torsion products, relative to the module M and to the exact sequence \(\mathcal{E}\) , the composite homomorphism

\[
\partial (\mathrm{M}, \mathcal {E}) = \psi_ {\mathrm{M}} \left(\mathrm{N} ^ {\prime}\right) ^ {- 1} \circ \partial \left(^ {\mathrm{M}} \mathcal {E}\right) \circ \psi_ {\mathrm{M}} \left(\mathrm{N} ^ {\prime \prime}\right): \operatorname{Tor} ^ {\mathrm{A}} \left(\mathrm{M}, \mathrm{N} ^ {\prime \prime}\right)\rightarrow \operatorname{Tor} ^ {\mathrm{A}} \left(\mathrm{M}, \mathrm{N} ^ {\prime}\right).
\]

This is a graded \(k\) -homomorphism of degree \((-1)\) , whose homogeneous components are denoted \(\partial_n(\mathbf{M}, \mathcal{E}) : \mathrm{Tor}_n^{\mathbf{A}}(\mathbf{M}, \mathbf{N}'') \to \mathrm{Tor}_{n-1}^{\mathbf{A}}(\mathbf{M}, \mathbf{N}')\) .

THEOREM 1. --- The left-unbounded sequence of homomorphisms of k-modules

\[
\begin{array}{l} \dots \longrightarrow \operatorname{Tor} _ {n} ^ {\mathbf {A}} (\mathbf {M}, \mathbf {N} ^ {\prime}) \xrightarrow {\operatorname{Tor} _ {n} ^ {\mathbf {A}} (1 , u)} \operatorname{Tor} _ {n} ^ {\mathbf {A}} (\mathbf {M}, \mathbf {N}) \xrightarrow {\operatorname{Tor} _ {n} ^ {\mathbf {A}} (1 , v)} \operatorname{Tor} _ {n} ^ {\mathbf {A}} (\mathbf {M}, \mathbf {N} ^ {\prime \prime}) \\ \xrightarrow {\hat {c} _ {n} (\mathbf {M} , \mathcal {E})} \operatorname{Tor} _ {n - 1} ^ {\mathbf {A}} (\mathbf {M}, \mathbf {N} ^ {\prime}) \xrightarrow {\operatorname{Tor} _ {n - 1} ^ {\mathbf {A}} (1 , u)} \dots \xrightarrow {\operatorname{Tor} _ {1} ^ {\mathbf {A}} (1 , v)} \operatorname{Tor} _ {1} ^ {\mathbf {A}} (\mathbf {M}, \mathbf {N} ^ {\prime \prime}) \\ \xrightarrow {\hat {c} _ {1} (\mathbf {M} , \mathcal {E})} \mathbf {M} \otimes_ {\mathbf {A}} \mathbf {N} ^ {\prime} \xrightarrow {1 \otimes u} \mathbf {M} \otimes_ {\mathbf {A}} \mathbf {N} \xrightarrow {1 \otimes v} \mathbf {M} \otimes_ {\mathbf {A}} \mathbf {N} ^ {\prime \prime} \longrightarrow 0 \end{array}
\]

is exact.

Consider indeed the diagram

\[
\begin{array}{c} \operatorname{Tor} (\mathbf {M}, \mathbf {N} ^ {\prime}) \xrightarrow {\operatorname{Tor} (1 , u)} \operatorname{Tor} (\mathbf {M}, \mathbf {N}) \xrightarrow {\operatorname{Tor} (1 , v)} \operatorname{Tor} (\mathbf {M}, \mathbf {N} ^ {\prime \prime}) \xrightarrow {\partial (\mathbf {M} , \mathcal {E})} \operatorname{Tor} (\mathbf {M}, \mathbf {N} ^ {\prime}) \xrightarrow {\operatorname{Tor} (1 , u)} \operatorname{Tor} (\mathbf {M}, \mathbf {N}) \\ \psi_ {\mathbf {M} (\mathbf {N} ^ {\prime})} \Big | _ {\downarrow} \quad \psi_ {\mathbf {M} (\mathbf {N})} \Big | _ {\downarrow} \quad \psi_ {\mathbf {M} (\mathbf {N} ^ {\prime \prime})} \Big | _ {\downarrow} \quad \psi_ {\mathbf {M} (\mathbf {N} ^ {\prime})} \Big | _ {\downarrow} \quad \psi_ {\mathbf {M} (\mathbf {N} ^ {\prime})} \Big | _ {\downarrow} \\ H (L (\mathbf {M}) \otimes N ^ {\prime}) \xrightarrow {H (1 \otimes u)} H (L (\mathbf {M}) \otimes N) \xrightarrow {H (1 \otimes v)} H (L (\mathbf {M} \otimes N ^ {\prime \prime}) \xrightarrow {\partial^ {(M)} \mathcal {E}} H (L (\mathbf {M}) \otimes N ^ {\prime}) \xrightarrow {} H (L (\mathbf {M}) \otimes N). \end{array}
\]

It is commutative by (X, p.~69, remark 3) and def. 2. On the other hand, the bottom row is exact (X, p.~30, th. 1), and the various \(\psi_{\mathrm{M}}\) are bijective (X, p.~69, prop. 5).

COROLLARY 1. --- If Tor \(_{1}^{A}\) (M, N'') = 0, the sequence

\[
0 \longrightarrow \mathrm{M} \otimes_ {\mathrm{A}} \mathrm{N} ^ {\prime} \xrightarrow {1 \otimes u} \mathrm{M} \otimes_ {\mathrm{A}} \mathrm{N} \xrightarrow {1 \otimes v} \mathrm{M} \otimes_ {\mathrm{A}} \mathrm{N} ^ {\prime \prime} \longrightarrow 0
\]

is exact.

COROLLARY 2. --- Let \(0 \to C' \xrightarrow{u} C \xrightarrow{v} C'' \to 0\) be an exact sequence of complexes of left A-modules and E a complex of right A-modules. If \(C''\) or E is flat, the sequence

\[
0 \longrightarrow \mathrm{E} \otimes_ {\mathrm{A}} \mathrm{C} ^ {\prime} \xrightarrow {1 \otimes u} \mathrm{E} \otimes_ {\mathrm{A}} \mathrm{C} \xrightarrow {1 \otimes v} \mathrm{E} \otimes_ {\mathrm{A}} \mathrm{C} ^ {\prime \prime} \longrightarrow 0
\]

is exact.

Indeed, \(\operatorname{Tor}_1^{\mathbf{A}}(\mathbf{E},\mathbf{C}^{\prime \prime}) = 0\) according to X, p.~69, cor. to prop. 5.

Example. --- Let a be an ideal of A. The exact sequence

\[
0 \rightarrow \mathfrak {a} \rightarrow \mathrm{A} _ {\mathrm{s}} \rightarrow \mathrm{A} / \mathfrak {a} \rightarrow 0
\]

of left A-modules gives rise to an exact sequence of torsion products, in which the terms \(\operatorname{Tor}_{i}^{\mathbf{A}}(\mathbf{M},\mathbf{A})\) are zero for i \textgreater{} 0. From this, we deduce isomorphisms

\[
\operatorname{Tor} _ {i + 1} ^ {\mathbf {A}} (\mathbf {M}, \mathbf {A} / \mathfrak {a}) \rightarrow \operatorname{Tor} _ {i} ^ {\mathbf {A}} (\mathbf {M}, \mathfrak {a}), \quad i > 0
\]

and an exact sequence

\[
0 \rightarrow \operatorname{Tor} _ {1} ^ {\mathbf {A}} (\mathrm{M}, \mathrm{A} / \mathfrak {a}) \rightarrow \mathrm{M} \otimes_ {\mathbf {A}} \mathfrak {a} \rightarrow \mathrm{M} \otimes_ {\mathbf {A}} \mathrm{A} \rightarrow \mathrm{M} \otimes \mathrm{A} / \mathfrak {a} \rightarrow 0:
\]

it follows that \(\operatorname{Tor}_{1}^{\mathbf{A}}(\mathbf{M},\mathbf{A}/\mathfrak{a})\) is identified with the kernel of the canonical homomorphism \(\mathbf{M}\otimes_{\mathbf{A}}\mathfrak{a}\to\mathbf{M}\) .

For example, taking for M a module of the form \(A_{d}/b\) where b is a right ideal of A, one obtains an isomorphism of \(\operatorname{Tor}_{1}^{A}(A/b, A/a)\) onto \((a \cap b)/ba\) .

PROPOSITION 9. --- Let \(f: \mathbf{M} \to \mathbf{M}_{1}\) be a homomorphism of right A-modules and

(ε)

\[
\begin{array}{c} 0 \longrightarrow \mathrm{N} ^ {\prime} \stackrel {{u}} {{\longrightarrow}} \mathrm{N} \stackrel {{v}} {{\longrightarrow}} \mathrm{N} ^ {\prime \prime} \longrightarrow 0 \\ g ^ {\prime} \Biggl \downarrow \qquad \qquad \qquad \qquad \qquad \qquad \qquad \qquad \qquad \qquad \qquad \qquad \qquad \qquad \qquad \qquad \qquad \qquad \qquad \qquad \qquad \qquad \qquad \qquad \qquad \qquad \qquad \qquad \qquad \qquad \qquad \qquad \qquad \qquad \\ 0 \longrightarrow \mathrm{N} _ {1} ^ {\prime} \stackrel {{u _ {1}}} {{\longrightarrow}} \mathrm{N} _ {1} \stackrel {{v _ {1}}} {{\longrightarrow}} \mathrm{N} _ {1} ^ {\prime \prime} \longrightarrow 0 \end{array}\tag{$(\mathcal{E}_1)$}
\]

a commutative diagram with exact rows of homomorphisms of left A-modules. The diagram of k-modules

\[
\begin{array}{c} \operatorname{Tor} ^ {\mathbf {A}} (\mathbf {M}, \mathbf {N} ^ {\prime \prime}) \xrightarrow {\hat {c} (\mathbf {M} , \mathcal {E})} \operatorname{Tor} ^ {\mathbf {A}} (\mathbf {M}, \mathbf {N} ^ {\prime}) \\ \operatorname{Tor} ^ {\mathbf {A}} (f, g ^ {\prime \prime}) \Big \downarrow \\ \operatorname{Tor} ^ {\mathbf {A}} (\mathbf {M} _ {1}, \mathbf {N} _ {1} ^ {\prime \prime}) \xrightarrow {\hat {c} (\mathbf {M} _ {1} , \mathcal {E} _ {1})} \operatorname{Tor} ^ {\mathbf {A}} (\mathbf {M} _ {1}, \mathbf {N} _ {1} ^ {\prime}) \end{array}
\]

is commutative.

This follows from X, p.~31, prop. 2, applied to the commutative diagram

\[
\begin{array}{l} 0 \longrightarrow \mathrm{L(M)} \otimes_ {\mathrm{A}} \mathrm{N} ^ {\prime} \xrightarrow {1 \otimes u} \mathrm{L(M)} \otimes_ {\mathrm{A}} \mathrm{N} \xrightarrow {1 \otimes v} \mathrm{L(M)} \otimes_ {\mathrm{A}} \mathrm{N} ^ {\prime \prime} \longrightarrow 0 \\ \quad \mathrm{L(f)} \otimes g ^ {\prime} \Biggl \downarrow \quad \mathrm{L(f)} \otimes g \Biggl \downarrow \quad \mathrm{L(f)} \otimes g ^ {\prime \prime} \Biggl \downarrow \\ 0 \longrightarrow \mathrm{L(M} _ {1}) \otimes_ {\mathrm{A}} \mathrm{N} _ {1} ^ {\prime} \xrightarrow {1 \otimes u _ {1}} \mathrm{L(M} _ {1}) \otimes_ {\mathrm{A}} \mathrm{N} _ {1} \xrightarrow {1 \otimes v _ {1}} \mathrm{L(M} _ {1}) \otimes_ {\mathrm{A}} \mathrm{N} _ {1} ^ {\prime \prime} \longrightarrow 0. \end{array}
\]

Analogously, if N is a left A-module and

\[
0 \to \mathbf {M} ^ {\prime} \xrightarrow {r} \mathbf {M} \xrightarrow {s} \mathbf {M} ^ {\prime \prime} \to 0\tag{F}
\]

an exact sequence of right A-modules, one defines connecting homomorphisms

\[
\partial (\mathcal {F}, \mathrm{N}): \operatorname{Tor} ^ {\mathrm{A}} \left(\mathrm{M} ^ {\prime \prime}, \mathrm{N}\right)\rightarrow \operatorname{Tor} ^ {\mathrm{A}} \left(\mathrm{M} ^ {\prime}, \mathrm{N}\right)
\]

\[
\partial_ {n} (\mathcal {F}, \mathrm{N}): \operatorname{Tor} _ {n} ^ {\mathrm{A}} \left(\mathrm{M} ^ {\prime \prime}, \mathrm{N}\right)\rightarrow \operatorname{Tor} _ {n - 1} ^ {\mathrm{A}} \left(\mathrm{M} ^ {\prime}, \mathrm{N}\right)
\]

by \(\partial(\mathcal{F},\mathrm{N})=\overline{\psi}_{\mathrm{N}}(\mathrm{M}^{\prime})^{-1}\circ\partial(\mathcal{F}^{\mathrm{N}})\circ\overline{\psi}_{\mathrm{N}}(\mathrm{M}^{\prime\prime})\) , where \(\partial(\mathcal{F}^{\mathrm{N}})\) is the connecting homomorphism of the exact sequence

\[
(\mathcal {F} ^ {N})
\]

\[
0 \rightarrow \mathrm{M} ^ {\prime} \otimes_ {\mathrm{A}} \mathrm{L} (\mathrm{N}) \rightarrow \mathrm{M} \otimes_ {\mathrm{A}} \mathrm{L} (\mathrm{N}) \rightarrow \mathrm{M} ^ {\prime \prime} \otimes_ {\mathrm{A}} \mathrm{L} (\mathrm{N}) \rightarrow 0
\]

deduced from \(\mathcal{F}\) , and we have:

THEOREM 1 bis. --- The left-unbounded sequence of homomorphisms of k-modules

\[
\begin{array}{l}\rightarrow \operatorname{Tor} _ {n} ^ {\mathsf {A}} (\mathbf {M} ^ {\prime}, \mathbf {N}) \xrightarrow {\operatorname{Tor} _ {n} ^ {\mathsf {A}} (r , 1)} \operatorname{Tor} _ {n} ^ {\mathsf {A}} (\mathbf {M}, \mathbf {N}) \xrightarrow {\operatorname{Tor} _ {n} ^ {\mathsf {A}} (s , 1)} \operatorname{Tor} _ {n} ^ {\mathsf {A}} (\mathbf {M} ^ {\prime \prime}, \mathbf {N}) \xrightarrow {\partial_ {n} (\mathcal {F} , \mathbf {N})} \operatorname{Tor} _ {n - 1} ^ {\mathsf {A}} (\mathbf {M} ^ {\prime}, \mathbf {N})\\\dots \rightarrow \operatorname{Tor} _ {1} ^ {\mathsf {A}} (\mathbf {M} ^ {\prime \prime}, \mathbf {N}) \xrightarrow {\partial_ {1} (\mathcal {F} , \mathbf {N})} \mathbf {M} ^ {\prime} \otimes_ {\mathbf {A}} \mathbf {N} \xrightarrow {r \otimes 1} \mathbf {M} \otimes_ {\mathbf {A}} \mathbf {N} \xrightarrow {s \otimes 1} \mathbf {M} ^ {\prime \prime} \otimes_ {\mathbf {A}} \mathbf {N} \longrightarrow 0\end{array}
\]

is exact.

We leave it to the reader to state and prove the properties analogous to the corollaries of Thm. 1 and to Prop. 9. Moreover:

PROPOSITION 10. --- Let ( \(F^{\circ}\) ) denote the exact sequence of left A-modules

\[
0 \rightarrow \mathrm{M} ^ {\prime \circ} \xrightarrow {r} \mathrm{M} ^ {\circ} \xrightarrow {s} \mathrm{M} ^ {\prime \prime \circ} \rightarrow 0.
\]

The diagram

\[
\begin{array}{c} \operatorname{Tor} ^ {\mathbf {A}} (\mathbf {M} ^ {\prime \prime}, \mathbf {N}) \xrightarrow {\partial (\mathcal {F} , \mathbf {N})} \operatorname{Tor} ^ {\mathbf {A}} (\mathbf {M} ^ {\prime}, \mathbf {N}) \\ \sigma_ {\mathbf {M} ^ {\prime}, \mathbf {N}} \Biggl \downarrow \qquad \qquad \qquad \qquad \qquad \qquad \qquad \qquad \qquad \sigma_ {\mathbf {M} ^ {\prime}, \mathbf {N}} \Biggl \downarrow \\ \operatorname{Tor} ^ {\mathbf {A} ^ {\circ}} (\mathbf {N} ^ {\circ}, \mathbf {M} ^ {\prime \prime^ {\circ}}) \xrightarrow {\partial (\mathbf {N} ^ {\circ} , \mathcal {F} ^ {\circ})} \operatorname{Tor} ^ {\mathbf {A} ^ {\circ}} (\mathbf {N} ^ {\circ}, \mathbf {M} ^ {\prime \circ}) \end{array}
\]

is commutative.

Indeed, this follows from X, p.~31, prop. 2, applied to the commutative diagram

\[
\begin{array}{l} 0 \longrightarrow \mathrm{M} ^ {\prime} \otimes_ {\mathrm{A}} \mathrm{L} (\mathrm{N}) \xrightarrow {r \otimes 1} \mathrm{M} \otimes_ {\mathrm{A}} \mathrm{L} (\mathrm{N}) \xrightarrow {s \otimes 1} \mathrm{M} ^ {\prime \prime} \otimes_ {\mathrm{A}} \mathrm{L} (\mathrm{N}) \longrightarrow 0 \\ \sigma (\mathrm{M} ^ {\prime}, \mathrm{L} (\mathrm{N})) \Biggl \downarrow \quad \sigma (\mathrm{M}, \mathrm{L} (\mathrm{N})) \Biggl \downarrow \quad \sigma (\mathrm{M} ^ {\prime \prime}, \mathrm{L} (\mathrm{N})) \Biggl \downarrow \\ 0 \longrightarrow \mathrm{L} (\mathrm{N} ^ {\circ}) \otimes_ {\mathrm{A} ^ {\circ}} \mathrm{M} ^ {\prime \circ} \xrightarrow {1 \otimes r} \mathrm{L} (\mathrm{N} ^ {\circ}) \otimes_ {\mathrm{A} ^ {\circ}} \mathrm{M} ^ {\circ} \xrightarrow {1 \otimes s} \mathrm{L} (\mathrm{N} ^ {\circ}) \otimes_ {\mathrm{A} ^ {\circ}} \mathrm{M} ^ {\prime \prime \circ} \longrightarrow 0. \end{array}
\]

We shall see other commutation relations later (cf. X, p. 131, cor. 1).

\subsection{Flat modules and torsion products}

THEOREM 2. --- Let E be a right A-module. The following conditions are equivalent:

\begin{enumerate}
\def\labelenumi{(\roman{enumi})}
\item
  E is flat;
\item
  for every left A-module F, and every integer n \textgreater{} 0, one has
\end{enumerate}

\[
\operatorname{Tor} _ {n} ^ {\mathbf {A}} (\mathbf {E}, \mathbf {F}) = 0;
\]

\begin{enumerate}
\def\labelenumi{(\roman{enumi})}
\setcounter{enumi}{2}
\tightlist
\item
  for every monogenic left A-module of finite presentation F, one has
\end{enumerate}

\[
\operatorname{Tor} _ {1} ^ {\mathbf {A}} (\mathbf {E}, \mathbf {F}) = 0;
\]

\begin{enumerate}
\def\labelenumi{(\roman{enumi})}
\setcounter{enumi}{3}
\item
  for every finitely generated left ideal a of A, the canonical map E ⊗\_A a → E is injective;
\item
  for every exact sequence of right A-modules, of the form
\end{enumerate}

\[
0 \rightarrow \mathbf {G} \xrightarrow {v} \mathbf {H} \xrightarrow {w} \mathbf {E} \rightarrow 0,
\]

and every left A-module F, the sequence

\[
0 \longrightarrow \mathrm{G} \otimes_ {\mathrm{A}} \mathrm{F} \xrightarrow {v \otimes 1} \mathrm{H} \otimes_ {\mathrm{A}} \mathrm{F} \xrightarrow {w \otimes 1} \mathrm{E} \otimes_ {\mathrm{A}} \mathrm{F} \longrightarrow 0
\]

is exact.

\begin{enumerate}
\def\labelenumi{(\roman{enumi})}
\item
  \(\Rightarrow\) (ii): this is the cor. to prop. 5 of X, p. 69.
\item
  \(\Rightarrow\) (iii): this is trivial.
\item
  \(\Leftrightarrow\) (iv): every monogenic left A-module of finite presentation is isomorphic to a quotient A/a, where a is a finitely generated left ideal, so that (iii) is equivalent to (iv) by X, p. 72, example.
\item
  \(\Rightarrow\) (i): by X, p. 8, prop. 3, X, p. 72, th. 1, E is flat as soon as \(\operatorname{Tor}_{1}^{\mathbf{A}}(\mathbf{E},\mathbf{F}) = 0\) for every left A-module F. If (iii) is satisfied, this is so as soon as F is monogenic and of finite presentation. By X, p. 11, prop. 7, every A-module (resp. every monogenic A-module) is a filtered inductive limit of modules of finite presentation (resp. of monogenic modules of finite presentation); one therefore sees, by X, p. 70, prop. 8, that it suffices to prove that if \(\operatorname{Tor}_{1}^{\mathbf{A}}(\mathbf{E},\mathbf{F}) = 0\) when F is monogenic, then this is so as soon as F is finitely generated. Let us therefore reason by induction on the cardinality of a generating family \((f_{1},\ldots,f_{n})\) of F; the exact sequence
\end{enumerate}

\[
0 \rightarrow \mathrm{A} f _ {1} \rightarrow \mathrm{F} \rightarrow \mathrm{F} / \mathrm{A} f _ {1} \rightarrow 0
\]

gives rise to an exact sequence

\[
\operatorname{Tor} _ {1} ^ {\mathbf {A}} (\mathrm{E}, \mathrm{A} f _ {1}) \rightarrow \operatorname{Tor} _ {1} ^ {\mathbf {A}} (\mathrm{E}, \mathrm{F}) \rightarrow \operatorname{Tor} _ {1} ^ {\mathbf {A}} (\mathrm{E}, \mathrm{F} / \mathrm{A} f _ {1}),
\]

so that \(\operatorname{Tor}_{1}^{\mathbf{A}}(\mathbf{E},\mathbf{F})=0\) since \(\operatorname{Tor}_{1}^{\mathbf{A}}(\mathbf{E},\mathbf{A}f_{1})=0\) and \(\operatorname{Tor}_{1}^{\mathbf{A}}(\mathbf{E},\mathbf{F}/\mathbf{A}f_{1})=0\) by the induction hypothesis.

\begin{enumerate}
\def\labelenumi{(\roman{enumi})}
\item
  \(\Rightarrow\) (v): this is cor. 2 to th. 1 (X, p. 72).
\item
  \(\Rightarrow\) (iii): the exact sequence (X, p. 50)
\end{enumerate}

\[
0 \longrightarrow Z _ {0} (\mathrm{E}) \stackrel {i _ {\mathrm{E}}} {\longrightarrow} L _ {0} (\mathrm{E}) \stackrel {p _ {\mathrm{E}}} {\longrightarrow} \mathrm{E} \longrightarrow 0
\]

gives rise, for every left A-module F, to an exact sequence

\[
0 \longrightarrow \operatorname{Tor} _ {1} ^ {\mathrm{A}} (\mathrm{E}, \mathrm{F}) \longrightarrow Z _ {0} (\mathrm{E}) \otimes_ {\mathrm{A}} \mathrm{F} \xrightarrow {i _ {\mathrm{E}} \otimes 1} \mathrm{L} _ {0} (\mathrm{E}) \otimes_ {\mathrm{A}} \mathrm{F} \xrightarrow {p _ {\mathrm{E}} \otimes 1} \mathrm{E} \otimes_ {\mathrm{A}} \mathrm{F} \longrightarrow 0.
\]

If (v) is satisfied, one has \(\operatorname{Tor}_1^{\mathbf{A}}(\mathbf{E},\mathbf{F}) = 0\) whence (iii).

COROLLARY 1. --- Let \(0 \to E' \to E \to E'' \to 0\) be an exact sequence of right A-modules. Suppose that \(E''\) is flat. Then for E to be flat, it is necessary and sufficient that \(E'\) be flat.

Let F be a left A-module. Since \(\operatorname{Tor}_i^{\mathbf{A}}(\mathbf{E}^{\prime \prime},\mathbf{F}) = 0\) for \(i = 1,2\) (th. 2, (i) \(\Rightarrow\) (ii)), one has an exact sequence

\[
0 \to \mathrm{Tor} _ {1} ^ {\mathbf {A}} (\mathbf {E} ^ {\prime}, \mathbf {F}) \to \mathrm{Tor} _ {1} ^ {\mathbf {A}} (\mathbf {E}, \mathbf {F}) \to 0
\]

whence the assertion (th. 2, (i) ⇔ (iii)).

COROLLARY 2. --- Let \(0 \to E_n \to E_{n-1} \to \cdots \to E_1 \to 0\) be an exact sequence of right A-modules. If \(E_i\) is flat for \(i = 1, \ldots, n-1\) , then \(E_n\) is flat.
